% GENERATED by tools/xref.py from the .aux files of the other volumes; do not edit.
\xrlabel{chap:1}{{1}{}{Light dark matter}{Doc-Start}{}}
\xrlabel{chap:1@cref}{{[chapter][1][]1}{[1][1][]1}}
\xrlabel{sec:1.1}{{1.1}{}{The master equation of direct detection}{Doc-Start}{}}
\xrlabel{sec:1.1@cref}{{[section][1][1]1.1}{[1][1][]1}}
\xrlabel{sec:1.2}{{1.2}{}{The local halo}{Doc-Start}{}}
\xrlabel{sec:1.2@cref}{{[section][2][1]1.2}{[1][1][]1}}
\xrlabel{eq:1.2-fgal}{{1.3}{}{Density and velocity distribution}{Doc-Start}{}}
\xrlabel{eq:1.2-fgal@cref}{{[equation][3][1]1.3}{[1][1][]1}}
\xrlabel{eq:1.2-nesc}{{1.4}{}{Density and velocity distribution}{Doc-Start}{}}
\xrlabel{eq:1.2-nesc@cref}{{[equation][4][1]1.4}{[1][1][]1}}
\xrlabel{tab:1-shm}{{1.1}{}{Standard-halo-model parameters adopted in this book, with their sources. Vectors are given in galactic axes (toward the Galactic center, along Galactic rotation, toward the north Galactic pole). The last two rows follow from the others. The reference Earth speed is the adopted magnitude of the solar velocity $(0, 238, 0) + (11.1, 12.2, 7.3)$\,\si {km.s^{-1}}. The largest laboratory speed is the sum of the escape speed and the reference Earth speed}{Doc-Start}{}}
\xrlabel{tab:1-shm@cref}{{[table][1][1]1.1}{[1][1][]1}}
\xrlabel{eq:1.2-vsun}{{1.5}{}{Density and velocity distribution}{Doc-Start}{}}
\xrlabel{eq:1.2-vsun@cref}{{[equation][5][1]1.5}{[1][1][]1}}
\xrlabel{eq:1.2-vE}{{1.6}{}{The laboratory frame}{Doc-Start}{}}
\xrlabel{eq:1.2-vE@cref}{{[equation][6][1]1.6}{[1][1][]1}}
\xrlabel{eq:1.2-flab}{{1.8}{}{The laboratory frame}{Doc-Start}{}}
\xrlabel{eq:1.2-flab@cref}{{[equation][8][1]1.8}{[1][1][]1}}
\xrlabel{eq:1.2-vmax}{{1.9}{}{The laboratory frame}{Doc-Start}{}}
\xrlabel{eq:1.2-vmax@cref}{{[equation][9][1]1.9}{[1][1][]1}}
\xrlabel{eq:1.2-Fspeed}{{1.10}{}{The laboratory frame}{Doc-Start}{}}
\xrlabel{eq:1.2-Fspeed@cref}{{[equation][10][1]1.10}{[1][1][]1}}
\xrlabel{eq:1.2-vorb}{{1.12}{}{Annual modulation}{Doc-Start}{}}
\xrlabel{eq:1.2-vorb@cref}{{[equation][12][1]1.12}{[1][1][]1}}
\xrlabel{eq:1.2-annual}{{1.13}{}{Annual modulation}{Doc-Start}{}}
\xrlabel{eq:1.2-annual@cref}{{[equation][13][1]1.13}{[1][1][]1}}
\xrlabel{eq:1.2-rotation}{{1.14}{}{Daily modulation}{Doc-Start}{}}
\xrlabel{eq:1.2-rotation@cref}{{[equation][14][1]1.14}{[1][1][]1}}
\xrlabel{fig:1-daily}{{1.3}{}{Angle $\theta _{\mathrm {lab}}$ between the Earth's galactic-frame velocity and the local vertical over one sidereal day, as a function of local sidereal time, at latitudes $90^\circ $N, $45^\circ $N, $0^\circ $ and $45^\circ $S. The velocity direction is the direction of the solar velocity, right ascension \SI {20.9}{h} and declination $47.5^\circ $ (vertical line: its transit), from Ref.~\cite [Table~1]{baxter2021recommended} and the coordinate rotation of Ref.~\cite [Eq.~(B3)]{mccabe2014earth}. The dark-matter wind arrives from this direction}{Doc-Start}{}}
\xrlabel{fig:1-daily@cref}{{[figure][3][1]1.3}{[1][1][]1}}
\xrlabel{eq:1.2-local}{{1.17}{}{Daily modulation}{Doc-Start}{}}
\xrlabel{eq:1.2-local@cref}{{[equation][17][1]1.17}{[1][1][]1}}
\xrlabel{sec:1.3}{{1.3}{}{Kinematics of slow dark matter}{Doc-Start}{}}
\xrlabel{sec:1.3@cref}{{[section][3][1]1.3}{[1][1][]1}}
\xrlabel{eq:1.3-ERmax}{{1.20}{}{Energy transfer and minimum velocity}{Doc-Start}{}}
\xrlabel{eq:1.3-ERmax@cref}{{[equation][20][1]1.20}{[1][1][]1}}
\xrlabel{eq:1.3-omegaq}{{1.21}{}{Energy transfer and minimum velocity}{Doc-Start}{}}
\xrlabel{eq:1.3-omegaq@cref}{{[equation][21][1]1.21}{[1][1][]1}}
\xrlabel{eq:1.3-vmin}{{1.22}{}{Energy transfer and minimum velocity}{Doc-Start}{}}
\xrlabel{eq:1.3-vmin@cref}{{[equation][22][1]1.22}{[1][1][]1}}
\xrlabel{eq:1.3-omegamax}{{1.23}{}{Energy transfer and minimum velocity}{Doc-Start}{}}
\xrlabel{eq:1.3-omegamax@cref}{{[equation][23][1]1.23}{[1][1][]1}}
\xrlabel{eq:1.3-mmin}{{1.24}{}{Threshold mass}{Doc-Start}{}}
\xrlabel{eq:1.3-mmin@cref}{{[equation][24][1]1.24}{[1][1][]1}}
\xrlabel{eq:1.3-region}{{1.25}{}{Momentum and energy scales}{Doc-Start}{}}
\xrlabel{eq:1.3-region@cref}{{[equation][25][1]1.25}{[1][1][]1}}
\xrlabel{fig:1-kinematics}{{1.4}{}{Momentum transfer $q$ and energy transfer $\omega $ accessible to elastic dark matter of mass 0.1, 1, 10 and \SI {100}{MeV} (blue curves, light to dark; shaded below), from $\omega \le qv_{\mathrm {max}} - q^2/2\mchi $ with $v_{\mathrm {max}} = \SI {794.5}{km.s^{-1}}$. The dotted line $\omega = qv_{\mathrm {max}}$ bounds all masses. Overlaid for silicon are the low-temperature band gap, \SI {1.17}{eV} (red), and the optical-phonon energy at the zone center, \SI {64}{meV} (green). The orange line is an acoustic branch $\omega $ = sound speed $\times \,q$ for a sound speed of \SI {5}{km.s^{-1}}, drawn up to $q = |\vect {G}_{111}|$, and the dashed curve is the recoil $q^2/2M$ of a free silicon nucleus, with $M = \SI {26.16}{GeV}$ the mean atomic mass of natural silicon. Vertical lines: the electron momentum scale $q_0 = \alpha m_e = \SI {3.73}{keV}$ and the length $|\vect {G}_{111}| = \sqrt {3}\,2\pi /a \approx \SI {4.0}{keV}$ of the shortest reciprocal-lattice vector of silicon, whose lattice constant is $a = 5.431$\,\AA . Scales from Refs.~\cites [Eqs.~(50) and (121), Sect.~IV\,A]{kahn2022searches}[Eq.~(11)]{ramanathan2020ionization}[lattice parameter of silicon]{mohr2025codata}}{Doc-Start}{}}
\xrlabel{fig:1-kinematics@cref}{{[figure][4][1]1.4}{[1][1][]1}}
\xrlabel{tab:1-kinematics}{{1.2}{}{Kinematic extremes for elastic dark matter at $v_{\mathrm {max}} = \SI {794.5}{km.s^{-1}}$. Columns: the dark-matter mass; the largest momentum transfer to a target that absorbs any momentum; the largest energy transfer; the largest recoil energy of a free silicon nucleus ($M = \SI {26.16}{GeV}$, the mean atomic mass of natural silicon). The $\tfrac 12\mchi v_{\mathrm {max}}^2$ column is also the threshold energy for which $m_{\mathrm {min}} = \mchi $}{Doc-Start}{}}
\xrlabel{tab:1-kinematics@cref}{{[table][2][1]1.2}{[1][1][]1}}
\xrlabel{lim:1.3a}{{1.1}{}{Nuclear recoils}{Doc-Start}{}}
\xrlabel{lim:1.3a@cref}{{[tcb@cnt@limitbox][1][1]1.1}{[1][1][]1}}
\xrlabel{eq:1.3-eta}{{1.27}{}{The speed integral}{Doc-Start}{}}
\xrlabel{eq:1.3-eta@cref}{{[equation][27][1]1.27}{[1][1][]1}}
\xrlabel{sec:1.3-edge}{{1.3.5}{}{The kinematic edge}{Doc-Start}{}}
\xrlabel{sec:1.3-edge@cref}{{[subsection][5][1,3]1.3.5}{[1][1][]1}}
\xrlabel{eq:1.3-edge}{{1.28}{}{The kinematic edge}{Doc-Start}{}}
\xrlabel{eq:1.3-edge@cref}{{[equation][28][1]1.28}{[1][1][]1}}
\xrlabel{sec:1.4}{{1.4}{}{Light dark-matter models}{Doc-Start}{}}
\xrlabel{sec:1.4@cref}{{[section][4][1]1.4}{[1][1][]1}}
\xrlabel{eq:1.4-yukawa}{{1.30}{}{Mediators and the Yukawa potential}{Doc-Start}{}}
\xrlabel{eq:1.4-yukawa@cref}{{[equation][30][1]1.30}{[1][1][]1}}
\xrlabel{eq:1.4-Vq}{{1.31}{}{Mediators and the Yukawa potential}{Doc-Start}{}}
\xrlabel{eq:1.4-Vq@cref}{{[equation][31][1]1.31}{[1][1][]1}}
\xrlabel{eq:1.4-defs}{{1.32}{}{Reference cross-section and form factor}{Doc-Start}{}}
\xrlabel{eq:1.4-defs@cref}{{[equation][32][1]1.32}{[1][1][]1}}
\xrlabel{eq:1.4-Vsq}{{1.33}{}{Reference cross-section and form factor}{Doc-Start}{}}
\xrlabel{eq:1.4-Vsq@cref}{{[equation][33][1]1.33}{[1][1][]1}}
\xrlabel{eq:1.4-yukawa-sigma}{{1.34}{}{Reference cross-section and form factor}{Doc-Start}{}}
\xrlabel{eq:1.4-yukawa-sigma@cref}{{[equation][34][1]1.34}{[1][1][]1}}
\xrlabel{lim:1.4a}{{1.2}{}{Reference cross-section and form factor}{Doc-Start}{}}
\xrlabel{lim:1.4a@cref}{{[tcb@cnt@limitbox][2][1]1.2}{[1][1][]1}}
\xrlabel{eq:1.4-sigmaN}{{1.36}{}{Scattering on nuclei}{Doc-Start}{}}
\xrlabel{eq:1.4-sigmaN@cref}{{[equation][36][1]1.36}{[1][1][]1}}
\xrlabel{eq:1.4-alp}{{1.39}{}{Absorption of bosonic dark matter}{Doc-Start}{}}
\xrlabel{eq:1.4-alp@cref}{{[equation][39][1]1.39}{[1][1][]1}}
\xrlabel{sec:1.5}{{1.5}{}{Detector signals}{Doc-Start}{}}
\xrlabel{sec:1.5@cref}{{[section][5][1]1.5}{[1][1][]1}}
\xrlabel{sec:1.5-quanta}{{1.5.1}{}{Readout quanta}{Doc-Start}{}}
\xrlabel{sec:1.5-quanta@cref}{{[subsection][1][1,5]1.5.1}{[1][1][]1}}
\xrlabel{eq:1.5-step}{{1.42}{}{Charge}{Doc-Start}{}}
\xrlabel{eq:1.5-step@cref}{{[equation][42][1]1.42}{[1][1][]1}}
\xrlabel{eq:1.5-luke}{{1.43}{}{Charge}{Doc-Start}{}}
\xrlabel{eq:1.5-luke@cref}{{[equation][43][1]1.43}{[1][1][]1}}
\xrlabel{eq:1.5-rate}{{1.46}{}{Event rate}{Doc-Start}{}}
\xrlabel{eq:1.5-rate@cref}{{[equation][46][1]1.46}{[1][1][]1}}
\xrlabel{sec:1.5-limit}{{1.5.3}{}{Counting and the upper limit}{Doc-Start}{}}
\xrlabel{sec:1.5-limit@cref}{{[subsection][3][1,5]1.5.3}{[1][1][]1}}
\xrlabel{eq:1.5-poisson}{{1.50}{}{Counting and the upper limit}{Doc-Start}{}}
\xrlabel{eq:1.5-poisson@cref}{{[equation][50][1]1.50}{[1][1][]1}}
\xrlabel{eq:1.5-N90}{{1.51}{}{Counting and the upper limit}{Doc-Start}{}}
\xrlabel{eq:1.5-N90@cref}{{[equation][51][1]1.51}{[1][1][]1}}
\xrlabel{eq:1.5-limit}{{1.52}{}{Counting and the upper limit}{Doc-Start}{}}
\xrlabel{eq:1.5-limit@cref}{{[equation][52][1]1.52}{[1][1][]1}}
\xrlabel{eq:1.5-limscaling}{{1.53}{}{Exclusion plots}{Doc-Start}{}}
\xrlabel{eq:1.5-limscaling@cref}{{[equation][53][1]1.53}{[1][1][]1}}
\xrlabel{sec:1.5-inputs}{{1.5.5}{}{Inputs of a limit}{Doc-Start}{}}
\xrlabel{sec:1.5-inputs@cref}{{[subsection][5][1,5]1.5.5}{[1][1][]1}}
\xrlabel{sec:2.1}{{2.1}{}{Fermi's golden rule}{Doc-Start}{}}
\xrlabel{sec:2.1@cref}{{[section][1][2]2.1}{[1][1][]1}}
\xrlabel{sec:2.2}{{2.2}{}{Born approximation and form factors}{Doc-Start}{}}
\xrlabel{sec:2.2@cref}{{[section][2][2]2.2}{[1][1][]1}}
\xrlabel{eq:2.2-dsdq}{{2.15}{}{The Born cross-section}{Doc-Start}{}}
\xrlabel{eq:2.2-dsdq@cref}{{[equation][15][2]2.15}{[1][1][]1}}
\xrlabel{eq:2.2-ff}{{2.17}{}{Scattering off a bound electron}{Doc-Start}{}}
\xrlabel{eq:2.2-ff@cref}{{[equation][17][2]2.17}{[1][1][]1}}
\xrlabel{eq:2.2-fion}{{2.22}{}{The hydrogen atom}{Doc-Start}{}}
\xrlabel{eq:2.2-fion@cref}{{[equation][22][2]2.22}{[1][1][]1}}
\xrlabel{eq:2.2-ion}{{2.23}{}{The hydrogen atom}{Doc-Start}{}}
\xrlabel{eq:2.2-ion@cref}{{[equation][23][2]2.23}{[1][1][]1}}
\xrlabel{eq:2.2-fion-pw}{{2.24}{}{The hydrogen atom}{Doc-Start}{}}
\xrlabel{eq:2.2-fion-pw@cref}{{[equation][24][2]2.24}{[1][1][]1}}
\xrlabel{eq:2.2-fermi}{{2.25}{}{The hydrogen atom}{Doc-Start}{}}
\xrlabel{eq:2.2-fermi@cref}{{[equation][25][2]2.25}{[1][1][]1}}
\xrlabel{lim:2.2a}{{2.1}{}{The hydrogen atom}{Doc-Start}{}}
\xrlabel{lim:2.2a@cref}{{[tcb@cnt@limitbox][1][2]2.1}{[1][1][]1}}
\xrlabel{eq:2.2-FN}{{2.26}{}{Nuclei and the Helm form factor}{Doc-Start}{}}
\xrlabel{eq:2.2-FN@cref}{{[equation][26][2]2.26}{[1][1][]1}}
\xrlabel{eq:2.2-nucleus}{{2.27}{}{Nuclei and the Helm form factor}{Doc-Start}{}}
\xrlabel{eq:2.2-nucleus@cref}{{[equation][27][2]2.27}{[1][1][]1}}
\xrlabel{eq:2.2-helm}{{2.28}{}{Nuclei and the Helm form factor}{Doc-Start}{}}
\xrlabel{eq:2.2-helm@cref}{{[equation][28][2]2.28}{[1][1][]1}}
\xrlabel{eq:2.3-nq}{{2.38}{}{One-body operators and the density}{Doc-Start}{}}
\xrlabel{eq:2.3-nq@cref}{{[equation][38][2]2.38}{[1][1][]1}}
\xrlabel{eq:2.3-HC}{{2.48}{}{Coulomb Hamiltonian of a solid}{Doc-Start}{}}
\xrlabel{eq:2.3-HC@cref}{{[equation][48][2]2.48}{[1][1][]1}}
\xrlabel{sec:2.4}{{2.4}{}{Second-order processes}{Doc-Start}{}}
\xrlabel{sec:2.4@cref}{{[section][4][2]2.4}{[1][1][]1}}
\xrlabel{eq:2.4-c2}{{2.51}{}{Interaction picture and Dyson series}{Doc-Start}{}}
\xrlabel{eq:2.4-c2@cref}{{[equation][51][2]2.51}{[1][1][]1}}
\xrlabel{eq:2.4-golden2}{{2.52}{}{Second-order golden rule}{Doc-Start}{}}
\xrlabel{eq:2.4-golden2@cref}{{[equation][52][2]2.52}{[1][1][]1}}
\xrlabel{eq:2.4-LS}{{2.53}{}{Second-order golden rule}{Doc-Start}{}}
\xrlabel{eq:2.4-LS@cref}{{[equation][53][2]2.53}{[1][1][]1}}
\xrlabel{eq:2.4-orderings}{{2.54}{}{Time orderings}{Doc-Start}{}}
\xrlabel{eq:2.4-orderings@cref}{{[equation][54][2]2.54}{[1][1][]1}}
\xrlabel{eq:2.4-lorentz}{{2.56}{}{Resonance and finite width}{Doc-Start}{}}
\xrlabel{eq:2.4-lorentz@cref}{{[equation][56][2]2.56}{[1][1][]1}}
\xrlabel{eq:2.4-width}{{2.57}{}{Resonance and finite width}{Doc-Start}{}}
\xrlabel{eq:2.4-width@cref}{{[equation][57][2]2.57}{[1][1][]1}}
\xrlabel{eq:2.4-SP}{{2.58}{}{Principal value and finite part}{Doc-Start}{}}
\xrlabel{eq:2.4-SP@cref}{{[equation][58][2]2.58}{[1][1][]1}}
\xrlabel{eq:2.4-split}{{2.60}{}{Principal value and finite part}{Doc-Start}{}}
\xrlabel{eq:2.4-split@cref}{{[equation][60][2]2.60}{[1][1][]1}}
\xrlabel{sec:2.5}{{2.5}{}{Thermal states and detailed balance}{Doc-Start}{}}
\xrlabel{sec:2.5@cref}{{[section][5][2]2.5}{[1][1][]1}}
\xrlabel{eq:2.5-transum}{{2.65}{}{Thermal golden rule}{Doc-Start}{}}
\xrlabel{eq:2.5-transum@cref}{{[equation][65][2]2.65}{[1][1][]1}}
\xrlabel{eq:2.5-golden}{{2.67}{}{Thermal golden rule}{Doc-Start}{}}
\xrlabel{eq:2.5-golden@cref}{{[equation][67][2]2.67}{[1][1][]1}}
\xrlabel{eq:2.5-db}{{2.68}{}{Detailed balance}{Doc-Start}{}}
\xrlabel{eq:2.5-db@cref}{{[equation][68][2]2.68}{[1][1][]1}}
\xrlabel{chap:3}{{3}{}{Electrons in crystals}{Doc-Start}{}}
\xrlabel{chap:3@cref}{{[chapter][3][]3}{[1][1][]1}}
\xrlabel{sec:3.1}{{3.1}{}{Lattices and reciprocal space}{Doc-Start}{}}
\xrlabel{sec:3.1@cref}{{[section][1][3]3.1}{[1][1][]1}}
\xrlabel{eq:3.1-qsplit}{{3.16}{}{Momentum transfer in a crystal}{Doc-Start}{}}
\xrlabel{eq:3.1-qsplit@cref}{{[equation][16][3]3.16}{[1][1][]1}}
\xrlabel{sec:3.2}{{3.2}{}{Bloch states and band structure}{Doc-Start}{}}
\xrlabel{sec:3.2@cref}{{[section][2][3]3.2}{[1][1][]1}}
\xrlabel{eq:3.2-rho}{{3.34}{}{The density vertex}{Doc-Start}{}}
\xrlabel{eq:3.2-rho@cref}{{[equation][34][3]3.34}{[1][1][]1}}
\xrlabel{eq:3.2-rhopw}{{3.35}{}{The density vertex}{Doc-Start}{}}
\xrlabel{eq:3.2-rhopw@cref}{{[equation][35][3]3.35}{[1][1][]1}}
\xrlabel{eq:3.2-closure}{{3.36}{}{The density vertex}{Doc-Start}{}}
\xrlabel{eq:3.2-closure@cref}{{[equation][36][3]3.36}{[1][1][]1}}
\xrlabel{eq:3.2-ne}{{3.37}{}{The density vertex}{Doc-Start}{}}
\xrlabel{eq:3.2-ne@cref}{{[equation][37][3]3.37}{[1][1][]1}}
\xrlabel{eq:3.2-dHk}{{3.40}{}{Velocity matrix elements}{Doc-Start}{}}
\xrlabel{eq:3.2-dHk@cref}{{[equation][40][3]3.40}{[1][1][]1}}
\xrlabel{eq:3.2-optical}{{3.43}{}{Velocity matrix elements}{Doc-Start}{}}
\xrlabel{eq:3.2-optical@cref}{{[equation][43][3]3.43}{[1][1][]1}}
\xrlabel{sec:3.3}{{3.3}{}{Silicon, germanium and gallium arsenide}{Doc-Start}{}}
\xrlabel{sec:3.3@cref}{{[section][3][3]3.3}{[1][1][]1}}
\xrlabel{eq:3.3-kc}{{3.47}{}{Band edges and valleys}{Doc-Start}{}}
\xrlabel{eq:3.3-kc@cref}{{[equation][47][3]3.47}{[1][1][]1}}
\xrlabel{fig:3-si-bands}{{3.7}{}{Silicon. (a)~Band structure along $\mathrm {L}$--$\mathrm {\Gamma }$--$\mathrm {X}$--$\mathrm {W}$--$\mathrm {K}$--$\mathrm {\Gamma }$ (points of \cref {sec:3.1}), valence bands black and conduction bands blue, energies from the valence-band maximum. The bands are the Kohn--Sham bands of the plane-wave calculation of \cref {app:E.1} at the measured lattice constant: generalized-gradient functional of Perdew, Burke and Ernzerhof (PBE), SG15 norm-conserving pseudopotential, cutoff \SI {60}{Ry}, $12^3$ grid for the density, $a = \SI {5.431}{\angstrom }$, no spin--orbit coupling~\cite {giannozzi2017advanced,hamann2013optimized,schlipf2015optimization}. The computed bands put the conduction-band minimum at $k_{\mathrm {c}} = 0.843\,(2\pi /a)$ (measured $0.85$) and the indirect gap $\Eg $ at \SI {0.557}{eV} (measured \SI {1.17}{eV}), the gap problem of density-functional theory. The vertical gap at $\mathrm {\Gamma }$, marked $E_{\mathrm {dir}}$, is \SI {2.56}{eV} (measured \SI {3.4}{eV}) and is the smallest vertical gap in the zone; the vertical gap at $\mathrm {L}$ is \SI {2.70}{eV} (measured \SI {3.45}{eV}). (b)~Minimum pair energy $E_{\mathrm {min}}(\vect {q})$, \cref {eq:3.3-Emin}, for $\vect {q}$ along [100], [110] and [111], from the same calculation on a grid of $48^3$ crystal momenta, with the conduction bands shifted rigidly up by \SI {0.61}{eV} (a scissor, \cref {sec:3.4}) so that the indirect gap equals the measured \SI {1.17}{eV} (dash-dotted line); at $\vect {q} = 0$ the shifted bands give $E_{\mathrm {min}} = \SI {3.17}{eV}$. Dashed: the largest energy that dark matter of \SI {1}{MeV} can deposit, $qv_{\mathrm {max}} - q^2/2\mchi $; dotted: the bound $qv_{\mathrm {max}}$ for any mass, with $v_{\mathrm {max}} = 2.650\times 10^{-3}$. The vertical line marks $|\vect {k}_{\mathrm {c}}| = \SI {1.94}{keV}$. Along [100], $E_{\mathrm {min}}$ returns to its value at $\vect {q} = 0$ at $|\vect {G}_{200}| = \SI {4.57}{keV}$}{Doc-Start}{}}
\xrlabel{fig:3-si-bands@cref}{{[figure][7][3]3.7}{[1][1][]1}}
\xrlabel{tab:3-materials}{{3.1}{}{Band parameters of silicon, germanium and GaAs. Lattice constant $a$ at room temperature: silicon~\cite [lattice parameter of silicon]{mohr2025codata}, germanium~\cite [Sect.~II\,A]{rideau2006strained}, GaAs~\cite {vurgaftman2001band}. Gaps at low temperature: silicon~\cite [Eq.~(11)]{ramanathan2020ionization}, germanium~\cite [Table~VI]{rideau2006strained}, GaAs~\cite [Table~I]{lucero2012improved}; at \SI {300}{K}~\cite [Table~I]{lucero2012improved}. Minimum direct gaps $E_{\mathrm {dir}}$ and spin--orbit splittings $\Delta _{\mathrm {SO}}$ of silicon and germanium~\cite [Table~II]{rideau2006strained}. Conduction valleys: position and number; masses along ($m^*_\parallel $) and across ($m^*_\perp $) the valley axis, silicon~\cite {hensel1965cyclotron} as listed in Ref.~\cite [Table~VII]{rideau2006strained}, germanium~\cite [Table~VII]{rideau2006strained}. Heavy- and light-hole masses $m_{\mathrm {hh}}$, $m_{\mathrm {lh}}$ along [100] and [111] from \cref {eq:3.3-holes} with the Luttinger parameters $(\gamma ^{\mathrm {Lu}}_1, \gamma ^{\mathrm {Lu}}_2, \gamma ^{\mathrm {Lu}}_3) = (4.285, 0.339, 1.446)$ for silicon and $(13.25, 4.20, 5.56)$ for germanium~\cite [Table~VII]{rideau2006strained}, and $(6.98, 2.06, 2.93)$ for GaAs. All GaAs entries not credited otherwise are from Ref.~\cite {vurgaftman2001band}. Masses in units of $m_e$}{Doc-Start}{}}
\xrlabel{tab:3-materials@cref}{{[table][1][3]3.1}{[1][1][]1}}
\xrlabel{eq:3.3-HSO}{{3.53}{}{Spin--orbit coupling}{Doc-Start}{}}
\xrlabel{eq:3.3-HSO@cref}{{[equation][53][3]3.53}{[1][1][]1}}
\xrlabel{eq:3.3-holes}{{3.56}{}{Spin--orbit coupling}{Doc-Start}{}}
\xrlabel{eq:3.3-holes@cref}{{[equation][56][3]3.56}{[1][1][]1}}
\xrlabel{eq:3.3-Emin}{{3.57}{}{The minimum pair energy}{Doc-Start}{}}
\xrlabel{eq:3.3-Emin@cref}{{[equation][57][3]3.57}{[1][1][]1}}
\xrlabel{eq:3.3-Eminpar}{{3.58}{}{The minimum pair energy}{Doc-Start}{}}
\xrlabel{eq:3.3-Eminpar@cref}{{[equation][58][3]3.58}{[1][1][]1}}
\xrlabel{eq:3.3-allowed}{{3.59}{}{The minimum pair energy}{Doc-Start}{}}
\xrlabel{eq:3.3-allowed@cref}{{[equation][59][3]3.59}{[1][1][]1}}
\xrlabel{we:3-kc}{{3.5}{}{The minimum pair energy}{Doc-Start}{}}
\xrlabel{we:3-kc@cref}{{[workedexample][5][3]3.5}{[1][1][]1}}
\xrlabel{we:3-mmin}{{3.6}{}{The minimum pair energy}{Doc-Start}{}}
\xrlabel{we:3-mmin@cref}{{[workedexample][6][3]3.6}{[1][1][]1}}
\xrlabel{sec:3.4}{{3.4}{}{Many electrons and density-functional theory}{Doc-Start}{}}
\xrlabel{sec:3.4@cref}{{[section][4][3]3.4}{[1][1][]1}}
\xrlabel{we:3-mmin-gap}{{3.7}{}{The scissor shift}{Doc-Start}{}}
\xrlabel{we:3-mmin-gap@cref}{{[workedexample][7][3]3.7}{[1][1][]1}}
\xrlabel{sec:4.1}{{4.1}{}{The Kubo formula}{Doc-Start}{}}
\xrlabel{sec:4.1@cref}{{[section][1][4]4.1}{[1][1][]1}}
\xrlabel{eq:4.1-kubo}{{4.3}{}{Response to a weak field}{Doc-Start}{}}
\xrlabel{eq:4.1-kubo@cref}{{[equation][3][4]4.3}{[1][1][]1}}
\xrlabel{sec:4.2}{{4.2}{}{Correlations and the dynamic structure factor}{Doc-Start}{}}
\xrlabel{sec:4.2@cref}{{[section][2][4]4.2}{[1][1][]1}}
\xrlabel{eq:4.2-Sab}{{4.13}{}{The response matrix}{Doc-Start}{}}
\xrlabel{eq:4.2-Sab@cref}{{[equation][13][4]4.13}{[1][1][]1}}
\xrlabel{eq:4.2-fdt}{{4.18}{}{The fluctuation--dissipation theorem}{Doc-Start}{}}
\xrlabel{eq:4.2-fdt@cref}{{[equation][18][4]4.18}{[1][1][]1}}
\xrlabel{eq:4.2-positivity}{{4.20}{}{Positivity}{Doc-Start}{}}
\xrlabel{eq:4.2-positivity@cref}{{[equation][20][4]4.20}{[1][1][]1}}
\xrlabel{we:4-gas}{{4.3}{}{Independent particles}{Doc-Start}{}}
\xrlabel{we:4-gas@cref}{{[workedexample][3][4]4.3}{[1][1][]1}}
\xrlabel{fig:4-gas}{{4.3}{}{Dynamic structure factor $S^0(q,\omega )$ of the non-interacting electron gas at zero temperature, in units of $m_ek_{\mathrm {F}}$, against the momentum transfer $q$ in units of the Fermi wavevector $k_{\mathrm {F}}$ and the energy $\omega $ in units of the Fermi energy $E_{\mathrm {F}} = k_{\mathrm {F}}^2/2m_e$. (a) Gray scale: $S^0$ inside the particle--hole continuum, bounded above by $\omega _+ = qv_{\mathrm {F}} + q^2/2m_e$ and below by $\omega _- = q^2/2m_e - qv_{\mathrm {F}}$ for $q > 2k_{\mathrm {F}}$ (solid). Below the dotted curve $\omega = qv_{\mathrm {F}} - q^2/2m_e$, $S^0 = m_e^2\omega /\pi q$ grows linearly in $\omega $; above it $S^0$ is a parabola in $\omega $ centered on the free recoil $q^2/2m_e$ (dashed). Red: the energies $\omega \le qv_{\mathrm {max}}$ that dark matter with the maximum halo speed $v_{\mathrm {max}} = 2.65\times 10^{-3}$ can deposit, for the valence density of silicon, $\nd _e = \SI {2.00e23}{cm^{-3}}$ ($k_{\mathrm {F}} = \SI {3.57}{keV}$, $E_{\mathrm {F}} = \SI {12.5}{eV}$, $v_{\mathrm {max}} = 0.379\,v_{\mathrm {F}}$). (b) $S^0$ against $\omega $ at $q = 0.5$, 1.5 and $3\,k_{\mathrm {F}}$. At $q = 0.5\,k_{\mathrm {F}}$ the triangle peaks at $0.239\,m_ek_{\mathrm {F}}$ at $\omega = 0.75\,E_{\mathrm {F}}$; at $q = 3\,k_{\mathrm {F}}$ the parabola spans $\omega _- = 3\,E_{\mathrm {F}}$ to $\omega _+ = 15\,E_{\mathrm {F}}$. Every curve integrates to $2\pi \nd _e$ times the fraction of the Fermi sea that can take the momentum $q$, and its first moment is $\pi \nd _eq^2/m_e$}{Doc-Start}{}}
\xrlabel{fig:4-gas@cref}{{[figure][3][4]4.3}{[1][1][]1}}
\xrlabel{sec:4.3}{{4.3}{}{Causality and sum rules}{Doc-Start}{}}
\xrlabel{sec:4.3@cref}{{[section][3][4]4.3}{[1][1][]1}}
\xrlabel{eq:4.3-KK}{{4.24}{}{Kramers--Kronig relations}{Doc-Start}{}}
\xrlabel{eq:4.3-KK@cref}{{[equation][24][4]4.24}{[1][1][]1}}
\xrlabel{eq:4.3-KKhalf}{{4.25}{}{Kramers--Kronig relations}{Doc-Start}{}}
\xrlabel{eq:4.3-KKhalf@cref}{{[equation][25][4]4.25}{[1][1][]1}}
\xrlabel{eq:4.3-firstmoment}{{4.28}{}{Moments from commutators}{Doc-Start}{}}
\xrlabel{eq:4.3-firstmoment@cref}{{[equation][28][4]4.28}{[1][1][]1}}
\xrlabel{eq:4.3-continuity}{{4.29}{}{The f-sum rule}{Doc-Start}{}}
\xrlabel{eq:4.3-continuity@cref}{{[equation][29][4]4.29}{[1][1][]1}}
\xrlabel{eq:4.3-fsum}{{4.30}{}{The f-sum rule}{Doc-Start}{}}
\xrlabel{eq:4.3-fsum@cref}{{[equation][30][4]4.30}{[1][1][]1}}
\xrlabel{eq:4.3-fsumNL}{{4.32}{}{Nonlocal potentials and the f-sum rule}{Doc-Start}{}}
\xrlabel{eq:4.3-fsumNL@cref}{{[equation][32][4]4.32}{[1][1][]1}}
\xrlabel{eq:4.3-mv}{{4.33}{}{Nonlocal potentials and the f-sum rule}{Doc-Start}{}}
\xrlabel{eq:4.3-mv@cref}{{[equation][33][4]4.33}{[1][1][]1}}
\xrlabel{eq:4.3-species}{{4.34}{}{Nonlocal potentials and the f-sum rule}{Doc-Start}{}}
\xrlabel{eq:4.3-species@cref}{{[equation][34][4]4.34}{[1][1][]1}}
\xrlabel{eq:4.3-atomq}{{4.35}{}{Oscillator strengths}{Doc-Start}{}}
\xrlabel{eq:4.3-atomq@cref}{{[equation][35][4]4.35}{[1][1][]1}}
\xrlabel{eq:4.3-osc}{{4.36}{}{Oscillator strengths}{Doc-Start}{}}
\xrlabel{eq:4.3-osc@cref}{{[equation][36][4]4.36}{[1][1][]1}}
\xrlabel{eq:4.3-TRK}{{4.37}{}{Oscillator strengths}{Doc-Start}{}}
\xrlabel{eq:4.3-TRK@cref}{{[equation][37][4]4.37}{[1][1][]1}}
\xrlabel{we:4-trk}{{4.5}{}{Oscillator strengths}{Doc-Start}{}}
\xrlabel{we:4-trk@cref}{{[workedexample][5][4]4.5}{[1][1][]1}}
\xrlabel{we:4-sifsum}{{4.7}{}{Oscillator strengths}{Doc-Start}{}}
\xrlabel{we:4-sifsum@cref}{{[workedexample][7][4]4.7}{[1][1][]1}}
\xrlabel{sec:4.4}{{4.4}{}{Screening and the dielectric function}{Doc-Start}{}}
\xrlabel{sec:4.4@cref}{{[section][4][4]4.4}{[1][1][]1}}
\xrlabel{eq:4.4-proper}{{4.39}{}{Proper response and the dielectric function}{Doc-Start}{}}
\xrlabel{eq:4.4-proper@cref}{{[equation][39][4]4.39}{[1][1][]1}}
\xrlabel{eq:4.4-epsmac}{{4.44}{}{Local fields and the dielectric matrix}{Doc-Start}{}}
\xrlabel{eq:4.4-epsmac@cref}{{[equation][44][4]4.44}{[1][1][]1}}
\xrlabel{eq:4.4-SL}{{4.45}{}{Structure factor and loss function}{Doc-Start}{}}
\xrlabel{eq:4.4-SL@cref}{{[equation][45][4]4.45}{[1][1][]1}}
\xrlabel{eq:4.4-fsumL}{{4.46}{}{Sum rules for the loss function}{Doc-Start}{}}
\xrlabel{eq:4.4-fsumL@cref}{{[equation][46][4]4.46}{[1][1][]1}}
\xrlabel{eq:4.4-fsumeps}{{4.47}{}{Sum rules for the loss function}{Doc-Start}{}}
\xrlabel{eq:4.4-fsumeps@cref}{{[equation][47][4]4.47}{[1][1][]1}}
\xrlabel{fig:4-siloss}{{4.7}{}{Dielectric function and loss function of crystalline silicon at $q \to 0$ and \SI {300}{K}, from the optical constants of Franta and coworkers, a Kramers--Kronig-consistent dispersion model fitted to ellipsometric and spectrophotometric measurements and tabulated from \SI {4}{meV} to \SI {40}{eV}~\cite {franta2017temperature}. The data are optical, that is transverse; at $q \to 0$ the transverse and longitudinal functions coincide (\cref {eq:4.4-ET}), and $L = \varepsilon _2/|\varepsilon |^2$. (a) $\varepsilon _1$ and $\varepsilon _2$; at low energy $\varepsilon _1 = \varepsilon _\infty = 11.7$ (dotted). The interband peaks of $\varepsilon _2$ lie at 3.4 and \SI {4.2}{eV}, and $\varepsilon _1$ crosses zero upward at \SI {16.3}{eV}. (b) The loss function on a logarithmic scale (solid), with its peak of 4.5 at \SI {16.4}{eV}. Dashed: the single damped oscillator of \cref {we:4-oscillatorL} at $q = 0$, which has the same f-sum and the same $\varepsilon _\infty $, and a damping equal to the measured full width, \SI {3.66}{eV}. Dotted: the room-temperature gap, \SI {1.12}{eV}, to match the data. Below \SI {40}{eV} the tabulated spectrum collects 96 percent of the f-sum of the four valence electrons in $\int \omega \varepsilon _2\,\dd \omega $ and 92 percent in $\int \omega L\,\dd \omega $}{Doc-Start}{}}
\xrlabel{fig:4-siloss@cref}{{[figure][7][4]4.7}{[1][1][]1}}
\xrlabel{eq:4.4-inverse}{{4.48}{}{Sum rules for the loss function}{Doc-Start}{}}
\xrlabel{eq:4.4-inverse@cref}{{[equation][48][4]4.48}{[1][1][]1}}
\xrlabel{eq:4.4-sigma}{{4.50}{}{Optical conductivity and transverse response}{Doc-Start}{}}
\xrlabel{eq:4.4-sigma@cref}{{[equation][50][4]4.50}{[1][1][]1}}
\xrlabel{eq:4.4-ET}{{4.52}{}{Optical conductivity and transverse response}{Doc-Start}{}}
\xrlabel{eq:4.4-ET@cref}{{[equation][52][4]4.52}{[1][1][]1}}
\xrlabel{eq:4.4-TF}{{4.53}{}{Thomas--Fermi screening}{Doc-Start}{}}
\xrlabel{eq:4.4-TF@cref}{{[equation][53][4]4.53}{[1][1][]1}}
\xrlabel{eq:4.4-dEdx}{{4.54}{}{Dielectric stopping}{Doc-Start}{}}
\xrlabel{eq:4.4-dEdx@cref}{{[equation][54][4]4.54}{[1][1][]1}}
\xrlabel{we:4-oscillatorL}{{4.9}{}{Dielectric stopping}{Doc-Start}{}}
\xrlabel{we:4-oscillatorL@cref}{{[workedexample][9][4]4.9}{[1][1][]1}}
\xrlabel{sec:4.5}{{4.5}{}{The homogeneous electron gas}{Doc-Start}{}}
\xrlabel{sec:4.5@cref}{{[section][5][4]4.5}{[1][1][]1}}
\xrlabel{eq:4.5-lindhard}{{4.56}{}{The independent-particle bubble}{Doc-Start}{}}
\xrlabel{eq:4.5-lindhard@cref}{{[equation][56][4]4.56}{[1][1][]1}}
\xrlabel{eq:4.5-moments}{{4.58}{}{The independent-particle bubble}{Doc-Start}{}}
\xrlabel{eq:4.5-moments@cref}{{[equation][58][4]4.58}{[1][1][]1}}
\xrlabel{eq:4.5-rpa}{{4.59}{}{The random-phase approximation}{Doc-Start}{}}
\xrlabel{eq:4.5-rpa@cref}{{[equation][59][4]4.59}{[1][1][]1}}
\xrlabel{eq:4.5-dispersion}{{4.60}{}{The plasmon}{Doc-Start}{}}
\xrlabel{eq:4.5-dispersion@cref}{{[equation][60][4]4.60}{[1][1][]1}}
\xrlabel{eq:4.5-plasmonL}{{4.61}{}{The plasmon}{Doc-Start}{}}
\xrlabel{eq:4.5-plasmonL@cref}{{[equation][61][4]4.61}{[1][1][]1}}
\xrlabel{eq:4.5-mermin}{{4.62}{}{Collisions and the Mermin function}{Doc-Start}{}}
\xrlabel{eq:4.5-mermin@cref}{{[equation][62][4]4.62}{[1][1][]1}}
\xrlabel{we:4-siplasmon}{{4.10}{}{Collisions and the Mermin function}{Doc-Start}{}}
\xrlabel{we:4-siplasmon@cref}{{[workedexample][10][4]4.10}{[1][1][]1}}
\xrlabel{sec:4.6}{{4.6}{}{Local fields and exchange--correlation kernels}{Doc-Start}{}}
\xrlabel{sec:4.6@cref}{{[section][6][4]4.6}{[1][1][]1}}
\xrlabel{eq:4.6-R0}{{4.63}{}{The crystal bubble}{Doc-Start}{}}
\xrlabel{eq:4.6-R0@cref}{{[equation][63][4]4.63}{[1][1][]1}}
\xrlabel{eq:4.6-rpaLF}{{4.68}{}{The random-phase approximation in a crystal}{Doc-Start}{}}
\xrlabel{eq:4.6-rpaLF@cref}{{[equation][68][4]4.68}{[1][1][]1}}
\xrlabel{eq:4.6-tddft}{{4.70}{}{Time-dependent density-functional theory}{Doc-Start}{}}
\xrlabel{eq:4.6-tddft@cref}{{[equation][70][4]4.70}{[1][1][]1}}
\xrlabel{eq:4.6-alda}{{4.71}{}{Time-dependent density-functional theory}{Doc-Start}{}}
\xrlabel{eq:4.6-alda@cref}{{[equation][71][4]4.71}{[1][1][]1}}
\xrlabel{eq:4.6-compress}{{4.73}{}{Local-field factor of the electron gas}{Doc-Start}{}}
\xrlabel{eq:4.6-compress@cref}{{[equation][73][4]4.73}{[1][1][]1}}
\xrlabel{we:4-siepsinf}{{4.11}{}{Many-body spread as an error measure}{Doc-Start}{}}
\xrlabel{we:4-siepsinf@cref}{{[workedexample][11][4]4.11}{[1][1][]1}}
\xrlabel{sec:4.7}{{4.7}{}{Green's functions}{Doc-Start}{}}
\xrlabel{sec:4.7@cref}{{[section][7][4]4.7}{[1][1][]1}}
\xrlabel{eq:4.7-qp}{{4.83}{}{Quasiparticles}{Doc-Start}{}}
\xrlabel{eq:4.7-qp@cref}{{[equation][83][4]4.83}{[1][1][]1}}
\xrlabel{sec:4.8}{{4.8}{}{Feynman diagrams}{Doc-Start}{}}
\xrlabel{sec:4.8@cref}{{[section][8][4]4.8}{[1][1][]1}}
\xrlabel{eq:4.8-W}{{4.99}{}{The screened interaction and the \texorpdfstring {$GW$}{GW} self-energy}{Doc-Start}{}}
\xrlabel{eq:4.8-W@cref}{{[equation][99][4]4.99}{[1][1][]1}}
\xrlabel{eq:4.8-lifetime}{{4.102}{}{The screened interaction and the \texorpdfstring {$GW$}{GW} self-energy}{Doc-Start}{}}
\xrlabel{eq:4.8-lifetime@cref}{{[equation][102][4]4.102}{[1][1][]1}}
\xrlabel{sec:4.9}{{4.9}{}{Excitons and the Bethe--Salpeter equation}{Doc-Start}{}}
\xrlabel{sec:4.9@cref}{{[section][9][4]4.9}{[1][1][]1}}
\xrlabel{eq:4.9-bse}{{4.104}{}{The electron--hole Hamiltonian}{Doc-Start}{}}
\xrlabel{eq:4.9-bse@cref}{{[equation][104][4]4.104}{[1][1][]1}}
\xrlabel{eq:4.9-kernels}{{4.105}{}{The electron--hole Hamiltonian}{Doc-Start}{}}
\xrlabel{eq:4.9-kernels@cref}{{[equation][105][4]4.105}{[1][1][]1}}
\xrlabel{eq:4.9-envelope}{{4.106}{}{Wannier--Mott excitons}{Doc-Start}{}}
\xrlabel{eq:4.9-envelope@cref}{{[equation][106][4]4.106}{[1][1][]1}}
\xrlabel{eq:4.9-WM}{{4.107}{}{Wannier--Mott excitons}{Doc-Start}{}}
\xrlabel{eq:4.9-WM@cref}{{[equation][107][4]4.107}{[1][1][]1}}
\xrlabel{eq:4.9-dispersion}{{4.108}{}{Wannier--Mott excitons}{Doc-Start}{}}
\xrlabel{eq:4.9-dispersion@cref}{{[equation][108][4]4.108}{[1][1][]1}}
\xrlabel{eq:4.9-epsmac}{{4.109}{}{Excitonic absorption and scattering}{Doc-Start}{}}
\xrlabel{eq:4.9-epsmac@cref}{{[equation][109][4]4.109}{[1][1][]1}}
\xrlabel{eq:4.9-eps2}{{4.110}{}{Excitonic absorption and scattering}{Doc-Start}{}}
\xrlabel{eq:4.9-eps2@cref}{{[equation][110][4]4.110}{[1][1][]1}}
\xrlabel{eq:4.9-elliott}{{4.111}{}{Excitonic absorption and scattering}{Doc-Start}{}}
\xrlabel{eq:4.9-elliott@cref}{{[equation][111][4]4.111}{[1][1][]1}}
\xrlabel{fig:4-bse}{{4.17}{}{The electron--hole attraction in the optical spectrum of crystalline silicon at $\vect {q} \to 0$, in a model calculation. Band energies and states come from the local empirical pseudopotential of Cohen and Bergstresser~\cite {cohen1966band}, fitted to measured transition energies, as in \cref {fig:4-si-approx}. All curves come from the transition-basis Hamiltonian \cref {eq:4.9-bse} in the Tamm--Dancoff approximation. The basis has 4 valence and 3 conduction bands on a shifted $12^3$ grid of $\vect {k}$ points over the full zone (20{,}736 transitions); both kernels use the 59 shortest reciprocal vectors, and the spectra carry a Lorentzian broadening of \SI {0.15}{eV}. The direct kernel uses the diagonal of the model's static inverse dielectric matrix in the random-phase approximation with local fields, computed with 26 conduction bands; spectra by the Lanczos--Haydock method. Black: measured at \SI {300}{K}~\cite {franta2017temperature}. (a) $\varepsilon _2$: independent particles (dashed gray), the random-phase approximation with local fields, that is, exchange kernel only (blue), and the Bethe--Salpeter equation with both kernels (red). The attraction turns the shoulder near \SI {3.4}{eV} into a peak; the whole spectrum shifts down because the fitted bands already contain part of the excitonic shift. (b) Loss function $L$ on a logarithmic scale for the same three cases (dotted blue: no local fields, which has the independent-particle $\varepsilon _2$). Below about \SI {3}{eV} the model curves are the Lorentzian tails of transitions that start at \SI {3.13}{eV} on this grid, and lie above the measured $L$, which has no such tails}{Doc-Start}{}}
\xrlabel{fig:4-bse@cref}{{[figure][17][4]4.17}{[1][1][]1}}
\xrlabel{eq:4.9-S}{{4.112}{}{Excitonic absorption and scattering}{Doc-Start}{}}
\xrlabel{eq:4.9-S@cref}{{[equation][112][4]4.112}{[1][1][]1}}
\xrlabel{chap:5}{{5}{}{Lattice dynamics}{Doc-Start}{}}
\xrlabel{chap:5@cref}{{[chapter][5][]5}{[1][1][]1}}
\xrlabel{sec:5.1}{{5.1}{}{The Born--Oppenheimer approximation}{Doc-Start}{}}
\xrlabel{sec:5.1@cref}{{[section][1][5]5.1}{[1][1][]1}}
\xrlabel{eq:5.1-dcoup}{{5.11}{}{The Hellmann--Feynman force}{Doc-Start}{}}
\xrlabel{eq:5.1-dcoup@cref}{{[equation][11][5]5.11}{[1][1][]1}}
\xrlabel{sec:5.2}{{5.2}{}{The harmonic crystal}{Doc-Start}{}}
\xrlabel{sec:5.2@cref}{{[section][2][5]5.2}{[1][1][]1}}
\xrlabel{eq:5.2-ortho}{{5.17}{}{Force constants and the dynamical matrix}{Doc-Start}{}}
\xrlabel{eq:5.2-ortho@cref}{{[equation][17][5]5.17}{[1][1][]1}}
\xrlabel{eq:5.2-sound}{{5.20}{}{Sound speeds}{Doc-Start}{}}
\xrlabel{eq:5.2-sound@cref}{{[equation][20][5]5.20}{[1][1][]1}}
\xrlabel{eq:5.2-dosac}{{5.22}{}{The density of states}{Doc-Start}{}}
\xrlabel{eq:5.2-dosac@cref}{{[equation][22][5]5.22}{[1][1][]1}}
\xrlabel{eq:5.2-debye}{{5.23}{}{The Debye model}{Doc-Start}{}}
\xrlabel{eq:5.2-debye@cref}{{[equation][23][5]5.23}{[1][1][]1}}
\xrlabel{we:5-debye}{{5.3}{}{The Debye model}{Doc-Start}{}}
\xrlabel{we:5-debye@cref}{{[workedexample][3][5]5.3}{[1][1][]1}}
\xrlabel{sec:5.3}{{5.3}{}{Phonons}{Doc-Start}{}}
\xrlabel{sec:5.3@cref}{{[section][3][5]5.3}{[1][1][]1}}
\xrlabel{eq:5.3-u}{{5.32}{}{Phonon operators}{Doc-Start}{}}
\xrlabel{eq:5.3-u@cref}{{[equation][32][5]5.32}{[1][1][]1}}
\xrlabel{eq:5.3-me}{{5.34}{}{Emission and absorption}{Doc-Start}{}}
\xrlabel{eq:5.3-me@cref}{{[equation][34][5]5.34}{[1][1][]1}}
\xrlabel{eq:5.3-B}{{5.35}{}{Zero-point and thermal motion}{Doc-Start}{}}
\xrlabel{eq:5.3-B@cref}{{[equation][35][5]5.35}{[1][1][]1}}
\xrlabel{tab:5-phonons}{{5.1}{}{Phonon energies and root-mean-square displacements $\langle u_x^2\rangle ^{1/2}$ along one Cartesian axis, in \AA , of silicon, germanium and GaAs. Phonon energies: measured by inelastic neutron scattering, silicon at room temperature~\cite {dolling1963inelastic,nilsson1972study}, germanium at \SI {80}{K}~\cite {nilsson1971phonon} as tabulated in Ref.~\cite [Table~IV]{haxhijaj2026lattice}, and GaAs~\cite {strauch1990phonon} as tabulated in Ref.~\cite [Table~I]{dehghani2025influence}; for GaAs the transverse (TO) and longitudinal (LO) optical energies at $\mathrm {\Gamma }$. Debye temperatures from the room-temperature elastic constants and lattice constants with the angular average of \cref {eq:5.2-cbar}: silicon from \cref {we:5-debye}, germanium from $C_{11} = 131$, $C_{12} = 49$, $C_{44} = \SI {68}{GPa}$~\cite {mcskimin1963elastic} (as tabulated in Ref.~\cite [Table~III]{haxhijaj2026lattice}) and the room-temperature $a = \SI {5.658}{\angstrom }$ of \cref {tab:3-materials}. Debye columns: \cref {eq:5.3-debye} with the natural isotope masses. DFPT columns: \cref {eq:5.3-u2} at 0 and \SI {300}{K} with the PBE density of states of \cref {fig:5-si} for silicon, on a $48^3$ grid of wavevectors. Germanium has no DFPT entry because PBE leaves it nearly metallic at the measured lattice constant, with a direct gap of \SI {0.05}{eV} and the wrong band order, and metallic at the PBE lattice constant. Entries marked $\dagger $ are computed for GaAs with the same method (PBE, $a = \SI {5.638}{\angstrom }$, a $4^3$ grid of wavevectors and the dipole term of \cref {sec:5.5}): the transverse acoustic energy at X, and the displacements of gallium (first) and arsenic (second) from the atom-projected densities of states $g_{\mathrm {ph},\kappa }$ on a $32^3$ grid. The germanium energies are tabulated to \SI {0.1}{THz}, \SI {0.4}{meV}. A dash marks a value not given here}{Doc-Start}{}}
\xrlabel{tab:5-phonons@cref}{{[table][1][5]5.1}{[1][1][]1}}
\xrlabel{sec:5.5}{{5.4}{}{Polar crystals}{Doc-Start}{}}
\xrlabel{sec:5.5@cref}{{[section][4][5]5.4}{[1][1][]1}}
\xrlabel{eq:5.5-U}{{5.38}{}{Born effective charges}{Doc-Start}{}}
\xrlabel{eq:5.5-U@cref}{{[equation][38][5]5.38}{[1][1][]1}}
\xrlabel{eq:5.5-Z}{{5.40}{}{Born effective charges}{Doc-Start}{}}
\xrlabel{eq:5.5-Z@cref}{{[equation][40][5]5.40}{[1][1][]1}}
\xrlabel{eq:5.5-neutral}{{5.41}{}{Born effective charges}{Doc-Start}{}}
\xrlabel{eq:5.5-neutral@cref}{{[equation][41][5]5.41}{[1][1][]1}}
\xrlabel{eq:5.5-na}{{5.42}{}{LO--TO splitting}{Doc-Start}{}}
\xrlabel{eq:5.5-na@cref}{{[equation][42][5]5.42}{[1][1][]1}}
\xrlabel{eq:5.5-LOTO}{{5.43}{}{LO--TO splitting}{Doc-Start}{}}
\xrlabel{eq:5.5-LOTO@cref}{{[equation][43][5]5.43}{[1][1][]1}}
\xrlabel{fig:5-gaas}{{5.4}{}{Phonon dispersion of GaAs. (a) Computed along $\mathrm {\Gamma }$--X--K--$\mathrm {\Gamma }$--L with density-functional perturbation theory: PBE functional, SG15 norm-conserving pseudopotentials with the gallium $3d$ shell in the valence, a \SI {60}{Ry} cutoff, an $8^3$ grid of crystal momenta, the static-lattice constant $a = \SI {5.638}{\angstrom }$ of \cref {tab:E-functionals} (the measured value extrapolated to zero temperature with the zero-point expansion removed) and natural masses, on a $4^3$ grid of wavevectors. The force constants are Fourier interpolated with the dipole term of \cref {eq:5.5-na} subtracted and added back, using $\varepsilon _\infty = 13.7$ and $Z^* = 2.094$ computed on a $28^3$ grid of crystal momenta; acoustic sum rules imposed on the force constants and on the Born charges. Circles: measured TO and LO energies at $\mathrm {\Gamma }$, 8.13 and \SI {8.79}{THz}~\cite {strauch1990phonon}, as tabulated in Ref.~\cite [Table~I]{dehghani2025influence}. (b) Optical branches near $\mathrm {\Gamma }$ along $[001]$ (left) and $[111]$ (right), with $|\vect {Q}|$ in units of $2\pi /a$. Diamond: LO energy at $\mathrm {\Gamma }$ from \cref {eq:5.5-LOTO} with the computed TO energy and Born charge and the measured $\varepsilon _\infty = 10.88$~\cite {lockwood2005optical}}{Doc-Start}{}}
\xrlabel{fig:5-gaas@cref}{{[figure][4][5]5.4}{[1][1][]1}}
\xrlabel{eq:5.5-eps}{{5.44}{}{The Lyddane--Sachs--Teller relation}{Doc-Start}{}}
\xrlabel{eq:5.5-eps@cref}{{[equation][44][5]5.44}{[1][1][]1}}
\xrlabel{eq:5.5-LST}{{5.45}{}{The Lyddane--Sachs--Teller relation}{Doc-Start}{}}
\xrlabel{eq:5.5-LST@cref}{{[equation][45][5]5.45}{[1][1][]1}}
\xrlabel{eq:D.14-LSTgen}{{5.46}{}{The Lyddane--Sachs--Teller relation}{Doc-Start}{}}
\xrlabel{eq:D.14-LSTgen@cref}{{[equation][46][5]5.46}{[1][1][]1}}
\xrlabel{we:5-lst}{{5.5}{}{The Lyddane--Sachs--Teller relation}{Doc-Start}{}}
\xrlabel{we:5-lst@cref}{{[workedexample][5][5]5.5}{[1][1][]1}}
\xrlabel{eq:5.5-HE}{{5.47}{}{Coupling to a uniform field}{Doc-Start}{}}
\xrlabel{eq:5.5-HE@cref}{{[equation][47][5]5.47}{[1][1][]1}}
\xrlabel{sec:5.6}{{5.5}{}{Displacement correlations}{Doc-Start}{}}
\xrlabel{sec:5.6@cref}{{[section][5][5]5.5}{[1][1][]1}}
\xrlabel{eq:5.6-W}{{5.48}{}{The Bloch identity}{Doc-Start}{}}
\xrlabel{eq:5.6-W@cref}{{[equation][48][5]5.48}{[1][1][]1}}
\xrlabel{eq:5.6-C}{{5.50}{}{The displacement correlation}{Doc-Start}{}}
\xrlabel{eq:5.6-C@cref}{{[equation][50][5]5.50}{[1][1][]1}}
\xrlabel{eq:5.6-gen}{{5.51}{}{The displacement correlation}{Doc-Start}{}}
\xrlabel{eq:5.6-gen@cref}{{[equation][51][5]5.51}{[1][1][]1}}
\xrlabel{eq:D.15-gen}{{5.52}{}{The displacement correlation}{Doc-Start}{}}
\xrlabel{eq:D.15-gen@cref}{{[equation][52][5]5.52}{[1][1][]1}}
\xrlabel{eq:5.6-multi}{{5.53}{}{Expansion in phonon number}{Doc-Start}{}}
\xrlabel{eq:5.6-multi@cref}{{[equation][53][5]5.53}{[1][1][]1}}
\xrlabel{eq:5.6-poisson}{{5.54}{}{Expansion in phonon number}{Doc-Start}{}}
\xrlabel{eq:5.6-poisson@cref}{{[equation][54][5]5.54}{[1][1][]1}}
\xrlabel{eq:5.6-sum}{{5.55}{}{Expansion in phonon number}{Doc-Start}{}}
\xrlabel{eq:5.6-sum@cref}{{[equation][55][5]5.55}{[1][1][]1}}
\xrlabel{we:5-dw}{{5.6}{}{Expansion in phonon number}{Doc-Start}{}}
\xrlabel{we:5-dw@cref}{{[workedexample][6][5]5.6}{[1][1][]1}}
\xrlabel{eq:5.6-FC}{{5.57}{}{The displaced oscillator}{Doc-Start}{}}
\xrlabel{eq:5.6-FC@cref}{{[equation][57][5]5.57}{[1][1][]1}}
\xrlabel{sec:5.7}{{5.6}{}{Anharmonicity}{Doc-Start}{}}
\xrlabel{sec:5.7@cref}{{[section][6][5]5.6}{[1][1][]1}}
\xrlabel{eq:5.7-asr}{{5.61}{}{Cubic and quartic force constants}{Doc-Start}{}}
\xrlabel{eq:5.7-asr@cref}{{[equation][61][5]5.61}{[1][1][]1}}
\xrlabel{eq:5.7-H3}{{5.62}{}{Cubic and quartic force constants}{Doc-Start}{}}
\xrlabel{eq:5.7-H3@cref}{{[equation][62][5]5.62}{[1][1][]1}}
\xrlabel{eq:5.7-V3}{{5.63}{}{Cubic and quartic force constants}{Doc-Start}{}}
\xrlabel{eq:5.7-V3@cref}{{[equation][63][5]5.63}{[1][1][]1}}
\xrlabel{eq:5.7-gamma}{{5.64}{}{Three-phonon decay and the linewidth}{Doc-Start}{}}
\xrlabel{eq:5.7-gamma@cref}{{[equation][64][5]5.64}{[1][1][]1}}
\xrlabel{we:5-klemens}{{5.8}{}{Three-phonon decay and the linewidth}{Doc-Start}{}}
\xrlabel{we:5-klemens@cref}{{[workedexample][8][5]5.8}{[1][1][]1}}
\xrlabel{eq:5.7-eta}{{5.66}{}{The anharmonic small parameter}{Doc-Start}{}}
\xrlabel{eq:5.7-eta@cref}{{[equation][66][5]5.66}{[1][1][]1}}
\xrlabel{eq:5.7-order}{{5.67}{}{The anharmonic small parameter}{Doc-Start}{}}
\xrlabel{eq:5.7-order@cref}{{[equation][67][5]5.67}{[1][1][]1}}
\xrlabel{eq:5.7-scha}{{5.69}{}{Effective harmonic theories}{Doc-Start}{}}
\xrlabel{eq:5.7-scha@cref}{{[equation][69][5]5.69}{[1][1][]1}}
\xrlabel{chap:6}{{6}{}{Electron--phonon coupling}{Doc-Start}{}}
\xrlabel{chap:6@cref}{{[chapter][6][]6}{[1][1][]1}}
\xrlabel{sec:6.1}{{6.1}{}{Nonadiabatic coupling}{Doc-Start}{}}
\xrlabel{sec:6.1@cref}{{[section][1][6]6.1}{[1][1][]1}}
\xrlabel{eq:6.1-sudden1}{{6.6}{}{Adiabatic and sudden limits}{Doc-Start}{}}
\xrlabel{eq:6.1-sudden1@cref}{{[equation][6][6]6.6}{[1][1][]1}}
\xrlabel{fig:6-sudden}{{6.2}{}{Adiabatic and sudden responses of an electron to a moving nucleus, in a model in which one electron is bound to the nucleus by a harmonic force of frequency $\omega _{\mathrm {e}}$ and the nucleus follows a prescribed classical path. (a) The velocity of the nucleus, $\dot X = \tfrac 12v_N[1 + \tanh (t/\tau _{\mathrm {acc}})]$, for acceleration times $\omega _{\mathrm {e}}\tau _{\mathrm {acc}} = 0$ (an instantaneous kick), 1 and 4. (b) The mean position $\langle y\rangle $ of the electron relative to the nucleus, in units of $v_N/\omega _{\mathrm {e}}$, for the same three cases. After a kick the electron oscillates about the nucleus with the full amplitude; after a slow acceleration it lags behind while the force acts and returns to rest. (c) The excitation $|\alpha _{\mathrm {d}}|^2 = -\ln P_0$, with $P_0$ the probability that the electron ends in the ground state, divided by its value $|\alpha _{\mathrm {d}}|^2_0 = m_ev_N^2/2\omega _{\mathrm {e}}$ for an instantaneous kick, against $\omega _{\mathrm {e}}\tau _{\mathrm {acc}}$. Solid: the exact result \cref {eq:6.1-model}. Dashed: its adiabatic asymptote $\pi ^2(\omega _{\mathrm {e}}\tau _{\mathrm {acc}})^2\ee ^{-\pi \omega _{\mathrm {e}}\tau _{\mathrm {acc}}}$. Dots: $-\ln P_0$ from a split-operator solution of the time-dependent Schr\"odinger equation for $m_ev_N^2/2\omega _{\mathrm {e}} = 1/2$. Shaded: the sudden region $\omega _{\mathrm {e}}\tau _{\mathrm {acc}} < 0.3$, where the excitation is within \SI {10}{\percent } of its sudden value}{Doc-Start}{}}
\xrlabel{fig:6-sudden@cref}{{[figure][2][6]6.2}{[1][1][]1}}
\xrlabel{eq:6.1-boost}{{6.8}{}{The sudden kick}{Doc-Start}{}}
\xrlabel{eq:6.1-boost@cref}{{[equation][8][6]6.8}{[1][1][]1}}
\xrlabel{eq:6.1-P0}{{6.9}{}{The sudden kick}{Doc-Start}{}}
\xrlabel{eq:6.1-P0@cref}{{[equation][9][6]6.9}{[1][1][]1}}
\xrlabel{sec:6.2}{{6.2}{}{Electron--phonon matrix elements}{Doc-Start}{}}
\xrlabel{sec:6.2@cref}{{[section][2][6]6.2}{[1][1][]1}}
\xrlabel{eq:6.2-screen}{{6.10}{}{The screened vertex}{Doc-Start}{}}
\xrlabel{eq:6.2-screen@cref}{{[equation][10][6]6.10}{[1][1][]1}}
\xrlabel{eq:6.2-g}{{6.11}{}{The screened vertex}{Doc-Start}{}}
\xrlabel{eq:6.2-g@cref}{{[equation][11][6]6.11}{[1][1][]1}}
\xrlabel{eq:6.2-asr}{{6.12}{}{Acoustic phonons and deformation potentials}{Doc-Start}{}}
\xrlabel{eq:6.2-asr@cref}{{[equation][12][6]6.12}{[1][1][]1}}
\xrlabel{eq:6.2-D}{{6.13}{}{Acoustic phonons and deformation potentials}{Doc-Start}{}}
\xrlabel{eq:6.2-D@cref}{{[equation][13][6]6.13}{[1][1][]1}}
\xrlabel{eq:6.2-frohlich}{{6.15}{}{The Fr\"ohlich vertex}{Doc-Start}{}}
\xrlabel{eq:6.2-frohlich@cref}{{[equation][15][6]6.15}{[1][1][]1}}
\xrlabel{sec:6.3}{{6.3}{}{Electron self-energy from phonons}{Doc-Start}{}}
\xrlabel{sec:6.3@cref}{{[section][3][6]6.3}{[1][1][]1}}
\xrlabel{eq:6.3-FM}{{6.19}{}{Fan--Migdal and Allen--Heine terms}{Doc-Start}{}}
\xrlabel{eq:6.3-FM@cref}{{[equation][19][6]6.19}{[1][1][]1}}
\xrlabel{eq:6.3-AH}{{6.20}{}{Fan--Migdal and Allen--Heine terms}{Doc-Start}{}}
\xrlabel{eq:6.3-AH@cref}{{[equation][20][6]6.20}{[1][1][]1}}
\xrlabel{eq:6.3-ren}{{6.21}{}{Fan--Migdal and Allen--Heine terms}{Doc-Start}{}}
\xrlabel{eq:6.3-ren@cref}{{[equation][21][6]6.21}{[1][1][]1}}
\xrlabel{eq:6.3-adiabatic}{{6.22}{}{The adiabatic limit}{Doc-Start}{}}
\xrlabel{eq:6.3-adiabatic@cref}{{[equation][22][6]6.22}{[1][1][]1}}
\xrlabel{we:6-zpr}{{6.3}{}{Zero-point renormalization}{Doc-Start}{}}
\xrlabel{we:6-zpr@cref}{{[workedexample][3][6]6.3}{[1][1][]1}}
\xrlabel{eq:6.3-ImFM}{{6.27}{}{Phonon-limited lifetimes}{Doc-Start}{}}
\xrlabel{eq:6.3-ImFM@cref}{{[equation][27][6]6.27}{[1][1][]1}}
\xrlabel{eq:6.3-cherenkov}{{6.28}{}{Phonon-limited lifetimes}{Doc-Start}{}}
\xrlabel{eq:6.3-cherenkov@cref}{{[equation][28][6]6.28}{[1][1][]1}}
\xrlabel{sec:6.4}{{6.4}{}{Phonon self-energy and the Eliashberg function}{Doc-Start}{}}
\xrlabel{sec:6.4@cref}{{[section][4][6]6.4}{[1][1][]1}}
\xrlabel{fig:6-bubble}{{6.6}{}{The two second-order diagrams of the electron--phonon interaction. Solid lines with arrows are electrons, wavy lines phonons, dots the vertex $g$. (a) The electron self-energy: an electron in $(n,\vect {k})$ exchanges a phonon $(\vect {Q},\nu )$ with the intermediate state $(m,\vect {k}+\vect {Q})$. (b) The phonon self-energy: a phonon creates an electron in $(m,\vect {k}+\vect {Q})$ and a hole in $(n,\vect {k})$ through $g$, and the pair recombines through $g^*$. (c) The Dyson series for the dressed phonon propagator $\mathsf {D}_\nu $ (double wavy line): the free propagator $\mathsf {D}_0$ with any number of bubbles inserted}{Doc-Start}{}}
\xrlabel{fig:6-bubble@cref}{{[figure][6][6]6.6}{[1][1][]1}}
\xrlabel{eq:6.4-Pi}{{6.29}{}{The electron bubble}{Doc-Start}{}}
\xrlabel{eq:6.4-Pi@cref}{{[equation][29][6]6.29}{[1][1][]1}}
\xrlabel{eq:6.4-pole}{{6.30}{}{The electron bubble}{Doc-Start}{}}
\xrlabel{eq:6.4-pole@cref}{{[equation][30][6]6.30}{[1][1][]1}}
\xrlabel{eq:6.4-allen}{{6.33}{}{The electron--phonon linewidth}{Doc-Start}{}}
\xrlabel{eq:6.4-allen@cref}{{[equation][33][6]6.33}{[1][1][]1}}
\xrlabel{eq:6.4-a2F}{{6.34}{}{The Eliashberg function}{Doc-Start}{}}
\xrlabel{eq:6.4-a2F@cref}{{[equation][34][6]6.34}{[1][1][]1}}
\xrlabel{eq:6.4-lambda}{{6.35}{}{The Eliashberg function}{Doc-Start}{}}
\xrlabel{eq:6.4-lambda@cref}{{[equation][35][6]6.35}{[1][1][]1}}
\xrlabel{eq:6.4-mass}{{6.36}{}{The Eliashberg function}{Doc-Start}{}}
\xrlabel{eq:6.4-mass@cref}{{[equation][36][6]6.36}{[1][1][]1}}
\xrlabel{fig:6-a2F}{{6.7}{}{Eliashberg function of aluminum (black line) in a model of free electrons scattered by screened empty-core ions, \cref {eq:6.4-a2Fgas}. The core radius, \SI {0.56}{\angstrom }, is taken where $\lambda _{\mathrm {ep}}$ is stationary with respect to it, and the phonons come from force constants between first and second neighbors fitted to measured zone-boundary frequencies and elastic constants. The gray area is the phonon density of states $g_{\mathrm {ph}}(\omega )$ of the same force constants, scaled to the height of the $\alpha ^2F$ peak. Dotted lines mark the transverse (TA) and longitudinal (LA) maxima of $g_{\mathrm {ph}}$. In $\alpha ^2F$ the transverse peak is suppressed, although three quarters of the modes lie below \SI {32}{meV}, because a transverse displacement couples to an electron only when the momentum transfer leaves the first zone. Both curves are histograms in bins of \SI {0.25}{meV} smoothed by a Gaussian of \SI {0.3}{meV}}{Doc-Start}{}}
\xrlabel{fig:6-a2F@cref}{{[figure][7][6]6.7}{[1][1][]1}}
\xrlabel{sec:6.5}{{6.5}{}{Phonon-assisted optical absorption}{Doc-Start}{}}
\xrlabel{sec:6.5@cref}{{[section][5][6]6.5}{[1][1][]1}}
\xrlabel{eq:6.5-HA}{{6.38}{}{The optical vertex}{Doc-Start}{}}
\xrlabel{eq:6.5-HA@cref}{{[equation][38][6]6.38}{[1][1][]1}}
\xrlabel{eq:6.5-eps2dir}{{6.40}{}{The optical vertex}{Doc-Start}{}}
\xrlabel{eq:6.5-eps2dir@cref}{{[equation][40][6]6.40}{[1][1][]1}}
\xrlabel{eq:6.5-M}{{6.42}{}{The phonon-assisted amplitude}{Doc-Start}{}}
\xrlabel{eq:6.5-M@cref}{{[equation][42][6]6.42}{[1][1][]1}}
\xrlabel{eq:6.5-eps2PA}{{6.43}{}{The phonon-assisted dielectric function}{Doc-Start}{}}
\xrlabel{eq:6.5-eps2PA@cref}{{[equation][43][6]6.43}{[1][1][]1}}
\xrlabel{eq:6.5-joint}{{6.44}{}{The Hall--Bardeen--Blatt edge}{Doc-Start}{}}
\xrlabel{eq:6.5-joint@cref}{{[equation][44][6]6.44}{[1][1][]1}}
\xrlabel{eq:6.5-edge}{{6.45}{}{The Hall--Bardeen--Blatt edge}{Doc-Start}{}}
\xrlabel{eq:6.5-edge@cref}{{[equation][45][6]6.45}{[1][1][]1}}
\xrlabel{eq:6.5-ratio}{{6.46}{}{Temperature dependence}{Doc-Start}{}}
\xrlabel{eq:6.5-ratio@cref}{{[equation][46][6]6.46}{[1][1][]1}}
\xrlabel{chap:7}{{7}{}{Magnets, superconductors, Dirac materials and defects}{Doc-Start}{}}
\xrlabel{chap:7@cref}{{[chapter][7][]7}{[1][1][]1}}
\xrlabel{sec:7.1}{{7.1}{}{Magnetic order and spin waves}{Doc-Start}{}}
\xrlabel{sec:7.1@cref}{{[section][1][7]7.1}{[1][1][]1}}
\xrlabel{eq:7.1-NiO}{{7.5}{}{Kinetic exchange}{Doc-Start}{}}
\xrlabel{eq:7.1-NiO@cref}{{[equation][5][7]7.5}{[1][1][]1}}
\xrlabel{eq:7.1-heis}{{7.6}{}{The Heisenberg model}{Doc-Start}{}}
\xrlabel{eq:7.1-heis@cref}{{[equation][6][7]7.6}{[1][1][]1}}
\xrlabel{eq:7.1-HP}{{7.7}{}{Spin waves of a ferromagnet}{Doc-Start}{}}
\xrlabel{eq:7.1-HP@cref}{{[equation][7][7]7.7}{[1][1][]1}}
\xrlabel{we:7-chain}{{7.1}{}{Spin waves of a ferromagnet}{Doc-Start}{}}
\xrlabel{we:7-chain@cref}{{[workedexample][1][7]7.1}{[1][1][]1}}
\xrlabel{eq:D.19-Jk}{{7.9}{}{Spin waves of an antiferromagnet}{Doc-Start}{}}
\xrlabel{eq:D.19-Jk@cref}{{[equation][9][7]7.9}{[1][1][]1}}
\xrlabel{eq:7.1-Hk}{{7.10}{}{Spin waves of an antiferromagnet}{Doc-Start}{}}
\xrlabel{eq:7.1-Hk@cref}{{[equation][10][7]7.10}{[1][1][]1}}
\xrlabel{eq:7.1-bogo}{{7.11}{}{Spin waves of an antiferromagnet}{Doc-Start}{}}
\xrlabel{eq:7.1-bogo@cref}{{[equation][11][7]7.11}{[1][1][]1}}
\xrlabel{eq:7.1-afm}{{7.12}{}{Spin waves of an antiferromagnet}{Doc-Start}{}}
\xrlabel{eq:7.1-afm@cref}{{[equation][12][7]7.12}{[1][1][]1}}
\xrlabel{we:7-nio}{{7.2}{}{Spin waves of an antiferromagnet}{Doc-Start}{}}
\xrlabel{we:7-nio@cref}{{[workedexample][2][7]7.2}{[1][1][]1}}
\xrlabel{fig:7-nio}{{7.1}{}{Magnetic order and spin waves of NiO. (a) Conventional cubic cell of the rock-salt structure: nickel ions (colored) on the fcc sites, oxygen ions (gray) between them. In the type-II order the spins of each $(111)$ plane are parallel and alternate in sign from plane to plane; two of these planes are shaded. The arrows show only the sign of each moment; in NiO the moments lie within the planes, along a direction set by the weak anisotropy. (b) Magnon energies $\omega ^{\mathrm {M}}_{\vect {k}}$ (meV) from linear spin-wave theory, \cref {eq:7.1-afm}, along three lines from the zone center $\Gamma $, with wavevectors in units of $2\pi /a$ and $a = \SI {4.17}{\angstrom }$. Model: undistorted fcc lattice, Heisenberg exchange with $J_1 = -\SI {1.37}{meV}$ and $J_2 = \SI {19.01}{meV}$ (each pair counted once, positive antiferromagnetic) measured by neutron scattering, $|\vect {S}| = 1$, no anisotropy. Solid: $J_1$ and $J_2$; dashed: $J_2$ alone, \cref {eq:7.1-J2}. Dotted: the linear laws $2\sqrt 3(J_1 + J_2)|\vect {S}|ak$ along $[111]$, and $2\sqrt 3J_2|\vect {S}|ak$ for $J_2$ alone. Dash-dotted: the measured maximum, ${\sim }\SI {117}{meV}$. The point $(\frac 12,\frac 12,\frac 12)$, along the stacking direction of the planes, is equivalent to $\Gamma $ in the magnetic structure, and the energy returns to zero there. The zeros at $(\frac 12,\frac 12,-\frac 12)$ and $(1,0,0)$ are soft points of the undistorted model; in real NiO the small rhombohedral distortion lifts them to a few meV. The largest energy in the zone, $|\vect {S}|(6J_2 - 3J_1) = \SI {118.2}{meV}$, lies off these lines}{Doc-Start}{}}
\xrlabel{fig:7-nio@cref}{{[figure][1][7]7.1}{[1][1][]1}}
\xrlabel{eq:7.1-Fmag}{{7.16}{}{Spin-density response}{Doc-Start}{}}
\xrlabel{eq:7.1-Fmag@cref}{{[equation][16][7]7.16}{[1][1][]1}}
\xrlabel{eq:7.1-Sfm}{{7.17}{}{Spin-density response}{Doc-Start}{}}
\xrlabel{eq:7.1-Sfm@cref}{{[equation][17][7]7.17}{[1][1][]1}}
\xrlabel{eq:D.19-inverse}{{7.18}{}{Spin-density response}{Doc-Start}{}}
\xrlabel{eq:D.19-inverse@cref}{{[equation][18][7]7.18}{[1][1][]1}}
\xrlabel{eq:7.1-weights}{{7.19}{}{Spin-density response}{Doc-Start}{}}
\xrlabel{eq:7.1-weights@cref}{{[equation][19][7]7.19}{[1][1][]1}}
\xrlabel{sec:7.2}{{7.2}{}{Superconductivity}{Doc-Start}{}}
\xrlabel{sec:7.2@cref}{{[section][2][7]7.2}{[1][1][]1}}
\xrlabel{eq:7.2-HBCS}{{7.23}{}{The reduced Hamiltonian}{Doc-Start}{}}
\xrlabel{eq:7.2-HBCS@cref}{{[equation][23][7]7.23}{[1][1][]1}}
\xrlabel{eq:7.2-gamma}{{7.25}{}{Mean field}{Doc-Start}{}}
\xrlabel{eq:7.2-gamma@cref}{{[equation][25][7]7.25}{[1][1][]1}}
\xrlabel{eq:7.2-E}{{7.26}{}{Mean field}{Doc-Start}{}}
\xrlabel{eq:7.2-E@cref}{{[equation][26][7]7.26}{[1][1][]1}}
\xrlabel{eq:7.2-dos}{{7.28}{}{Quasiparticle density of states}{Doc-Start}{}}
\xrlabel{eq:7.2-dos@cref}{{[equation][28][7]7.28}{[1][1][]1}}
\xrlabel{eq:7.2-gap}{{7.29}{}{The gap equation}{Doc-Start}{}}
\xrlabel{eq:7.2-gap@cref}{{[equation][29][7]7.29}{[1][1][]1}}
\xrlabel{eq:7.2-threshold}{{7.30}{}{The pair-breaking threshold}{Doc-Start}{}}
\xrlabel{eq:7.2-threshold@cref}{{[equation][30][7]7.30}{[1][1][]1}}
\xrlabel{eq:7.2-coh}{{7.31}{}{Coherence factors}{Doc-Start}{}}
\xrlabel{eq:7.2-coh@cref}{{[equation][31][7]7.31}{[1][1][]1}}
\xrlabel{eq:7.2-cohE}{{7.32}{}{Coherence factors}{Doc-Start}{}}
\xrlabel{eq:7.2-cohE@cref}{{[equation][32][7]7.32}{[1][1][]1}}
\xrlabel{eq:7.2-pairspec}{{7.33}{}{The pair-creation spectrum}{Doc-Start}{}}
\xrlabel{eq:7.2-pairspec@cref}{{[equation][33][7]7.33}{[1][1][]1}}
\xrlabel{eq:7.2-nambu}{{7.35}{}{Nambu form}{Doc-Start}{}}
\xrlabel{eq:7.2-nambu@cref}{{[equation][35][7]7.35}{[1][1][]1}}
\xrlabel{eq:7.2-Rc}{{7.37}{}{The conserving response}{Doc-Start}{}}
\xrlabel{eq:7.2-Rc@cref}{{[equation][37][7]7.37}{[1][1][]1}}
\xrlabel{eq:7.2-cancel}{{7.39}{}{The conserving response}{Doc-Start}{}}
\xrlabel{eq:7.2-cancel@cref}{{[equation][39][7]7.39}{[1][1][]1}}
\xrlabel{we:7-al}{{7.3}{}{The Ward identity}{Doc-Start}{}}
\xrlabel{we:7-al@cref}{{[workedexample][3][7]7.3}{[1][1][]1}}
\xrlabel{sec:7.3}{{7.3}{}{Dirac and narrow-gap materials}{Doc-Start}{}}
\xrlabel{sec:7.3@cref}{{[section][3][7]7.3}{[1][1][]1}}
\xrlabel{eq:7.3-H}{{7.42}{}{The Dirac Hamiltonian from \texorpdfstring {$\vect {k}\cdot \vect {p}$}{k.p} theory}{Doc-Start}{}}
\xrlabel{eq:7.3-H@cref}{{[equation][42][7]7.42}{[1][1][]1}}
\xrlabel{eq:7.3-edge}{{7.45}{}{Interband pair threshold}{Doc-Start}{}}
\xrlabel{eq:7.3-edge@cref}{{[equation][45][7]7.45}{[1][1][]1}}
\xrlabel{eq:7.3-epsbg}{{7.46}{}{Screening by Dirac electrons}{Doc-Start}{}}
\xrlabel{eq:7.3-epsbg@cref}{{[equation][46][7]7.46}{[1][1][]1}}
\xrlabel{eq:7.3-alpha}{{7.47}{}{Screening by Dirac electrons}{Doc-Start}{}}
\xrlabel{eq:7.3-alpha@cref}{{[equation][47][7]7.47}{[1][1][]1}}
\xrlabel{eq:7.3-eps}{{7.48}{}{Screening by Dirac electrons}{Doc-Start}{}}
\xrlabel{eq:7.3-eps@cref}{{[equation][48][7]7.48}{[1][1][]1}}
\xrlabel{eq:7.3-static}{{7.50}{}{Screening by Dirac electrons}{Doc-Start}{}}
\xrlabel{eq:7.3-static@cref}{{[equation][50][7]7.50}{[1][1][]1}}
\xrlabel{eq:D.22-gapless}{{7.51}{}{Screening by Dirac electrons}{Doc-Start}{}}
\xrlabel{eq:D.22-gapless@cref}{{[equation][51][7]7.51}{[1][1][]1}}
\xrlabel{eq:7.3-aniso}{{7.52}{}{Screening by Dirac electrons}{Doc-Start}{}}
\xrlabel{eq:7.3-aniso@cref}{{[equation][52][7]7.52}{[1][1][]1}}
\xrlabel{we:7-zrte5}{{7.4}{}{Screening by Dirac electrons}{Doc-Start}{}}
\xrlabel{we:7-zrte5@cref}{{[workedexample][4][7]7.4}{[1][1][]1}}
\xrlabel{sec:7.4}{{7.4}{}{Defects and dopants}{Doc-Start}{}}
\xrlabel{sec:7.4@cref}{{[section][4][7]7.4}{[1][1][]1}}
\xrlabel{eq:7.4-EMT}{{7.53}{}{Shallow donors}{Doc-Start}{}}
\xrlabel{eq:7.4-EMT@cref}{{[equation][53][7]7.53}{[1][1][]1}}
\xrlabel{fig:7-donor}{{7.5}{}{Bound levels of a shallow donor in silicon, measured from the conduction-band minimum (meV). Gray: the hydrogenic model, \cref {eq:7.4-hydrogenic}, with $\varepsilon _0 = 11.7$ and the transverse mass $0.191\,m_e$ ($1s$ at \SI {19.0}{meV}) or the longitudinal mass $0.916\,m_e$ ($1s$ at \SI {91.0}{meV}, off the scale), each with its $n = 2$ level. Blue: the effective-mass equation \cref {eq:7.4-EMT} with both masses, solved numerically (coupled radial equations in spherical harmonics after scaling $z$ by $\gamma _m^{1/2}$, $l \le 14$), for $\varepsilon _0 = 11.7$ (room temperature) and 11.4 (liquid-helium temperature); from the top, $3p_0$, $2p_\pm $, $2s$, $2p_0$ and $1s$. Orange: measured levels of phosphorus and arsenic donors~\cite [Fig.~1]{vinh2008silicon}. The $p$ levels agree with effective-mass theory at $\varepsilon _0 = 11.4$ to \SI {0.4}{\percent }; the $1s$ level splits into $A_1$, $T_2$ and $E$ states, and the $A_1$ ground state lies below the effective-mass value by \SI {14}{meV} for phosphorus and \SI {22}{meV} for arsenic}{Doc-Start}{}}
\xrlabel{fig:7-donor@cref}{{[figure][5][7]7.5}{[1][1][]1}}
\xrlabel{we:7-sip}{{7.5}{}{Shallow donors}{Doc-Start}{}}
\xrlabel{we:7-sip@cref}{{[workedexample][5][7]7.5}{[1][1][]1}}
\xrlabel{eq:7.4-Ed}{{7.59}{}{Displacement thresholds and Frenkel pairs}{Doc-Start}{}}
\xrlabel{eq:7.4-Ed@cref}{{[equation][59][7]7.59}{[1][1][]1}}
\xrlabel{chap:8}{{8}{}{The response framework}{Doc-Start}{}}
\xrlabel{chap:8@cref}{{[chapter][8][]8}{[1][1][]1}}
\xrlabel{sec:8.1}{{8.1}{}{Dark-matter couplings to a target}{Doc-Start}{}}
\xrlabel{sec:8.1@cref}{{[section][1][8]8.1}{[1][1][]1}}
\xrlabel{eq:8.1-rows}{{8.3}{}{Target operators per species}{Doc-Start}{}}
\xrlabel{eq:8.1-rows@cref}{{[equation][3][8]8.3}{[1][1][]1}}
\xrlabel{eq:8.1-O7O8}{{8.4}{}{Target operators per species}{Doc-Start}{}}
\xrlabel{eq:8.1-O7O8@cref}{{[equation][4][8]8.4}{[1][1][]1}}
\xrlabel{eq:8.1-transverse}{{8.5}{}{Target operators per species}{Doc-Start}{}}
\xrlabel{eq:8.1-transverse@cref}{{[equation][5][8]8.5}{[1][1][]1}}
\xrlabel{tab:8-operators}{{8.1}{}{Non-relativistic operators of dark-matter scattering and absorption, and the target operators each needs. Upper block: the operators $\mathcal {O}_i$ of Refs.~\cite [Eqs.~(19)--(24) and (45)]{fitzpatrick2013effective} and~\cite [Eqs.~(12) and (13)]{anand2014model}, written for an electron with $\vect {q}$ the momentum the dark matter loses; $\vect {v}^\perp $ is \cref {eq:8.1-vperp}, $\bar {\vect {v}} \equiv \bar {\vect {v}}_\chi = \vect {v} - \vect {q}/2\mchi $, and $m = m_e$. The rows are those of \cref {eq:8.1-rows} for $s = e$, all at the same $\vect {q}$; repeated Cartesian indices are summed, and $\epsilon _{\alpha \beta \gamma }$ is the Levi-Civita symbol. For nucleons, replace $e$ by $N$ and $m_e$ by the nucleon mass; the nuclear rows follow from \cref {eq:8.1-nuclear}. Lower block: absorption of a boson $\phi $ of mass $m_\phi $ at $\vect {q} = \vect {p}_\phi $, and fermionic absorption with an outgoing fermion $\tilde \nu $ (\cref {sec:8.1-absorption}). Last column: what each row excites in an insulator. Through the nuclei the density row also excites phonons, and in an ordered magnet the spin rows excite magnons (\cref {sec:7.1})}{Doc-Start}{}}
\xrlabel{tab:8-operators@cref}{{[table][1][8]8.1}{[1][1][]1}}
\xrlabel{eq:8.1-Hint}{{8.6}{}{Target operators per species}{Doc-Start}{}}
\xrlabel{eq:8.1-Hint@cref}{{[equation][6][8]8.6}{[1][1][]1}}
\xrlabel{eq:8.1-D}{{8.7}{}{Target operators per species}{Doc-Start}{}}
\xrlabel{eq:8.1-D@cref}{{[equation][7][8]8.7}{[1][1][]1}}
\xrlabel{sec:8.1-nuclei}{{8.1.3}{}{Nuclei}{Doc-Start}{}}
\xrlabel{sec:8.1-nuclei@cref}{{[subsection][3][8,1]8.1.3}{[1][1][]1}}
\xrlabel{eq:8.1-nuclear}{{8.8}{}{Nuclei}{Doc-Start}{}}
\xrlabel{eq:8.1-nuclear@cref}{{[equation][8][8]8.8}{[1][1][]1}}
\xrlabel{we:8-darkphoton}{{8.1}{}{Nuclei}{Doc-Start}{}}
\xrlabel{we:8-darkphoton@cref}{{[workedexample][1][8]8.1}{[1][1][]1}}
\xrlabel{sec:8.1-absorption}{{8.1.5}{}{Absorption operators}{Doc-Start}{}}
\xrlabel{sec:8.1-absorption@cref}{{[subsection][5][8,1]8.1.5}{[1][1][]1}}
\xrlabel{eq:8.1-alp}{{8.10}{}{Absorption operators}{Doc-Start}{}}
\xrlabel{eq:8.1-alp@cref}{{[equation][10][8]8.10}{[1][1][]1}}
\xrlabel{eq:8.1-scalar}{{8.11}{}{Absorption operators}{Doc-Start}{}}
\xrlabel{eq:8.1-scalar@cref}{{[equation][11][8]8.11}{[1][1][]1}}
\xrlabel{sec:8.1-kinematics}{{8.1.6}{}{Energy deposited}{Doc-Start}{}}
\xrlabel{sec:8.1-kinematics@cref}{{[subsection][6][8,1]8.1.6}{[1][1][]1}}
\xrlabel{eq:8.1-omega}{{8.12}{}{Energy deposited}{Doc-Start}{}}
\xrlabel{eq:8.1-omega@cref}{{[equation][12][8]8.12}{[1][1][]1}}
\xrlabel{eq:8.1-omega-inel}{{8.13}{}{Energy deposited}{Doc-Start}{}}
\xrlabel{eq:8.1-omega-inel@cref}{{[equation][13][8]8.13}{[1][1][]1}}
\xrlabel{eq:8.1-omega-ferm}{{8.14}{}{Energy deposited}{Doc-Start}{}}
\xrlabel{eq:8.1-omega-ferm@cref}{{[equation][14][8]8.14}{[1][1][]1}}
\xrlabel{sec:8.2}{{8.2}{}{The master rate formula}{Doc-Start}{}}
\xrlabel{sec:8.2@cref}{{[section][2][8]8.2}{[1][1][]1}}
\xrlabel{eq:8.2-rate}{{8.15}{}{Scattering}{Doc-Start}{}}
\xrlabel{eq:8.2-rate@cref}{{[equation][15][8]8.15}{[1][1][]1}}
\xrlabel{eq:8.2-spectrum}{{8.16}{}{Scattering}{Doc-Start}{}}
\xrlabel{eq:8.2-spectrum@cref}{{[equation][16][8]8.16}{[1][1][]1}}
\xrlabel{sec:8.2-conditions}{{8.2.2}{}{Assumptions of the formula}{Doc-Start}{}}
\xrlabel{sec:8.2-conditions@cref}{{[subsection][2][8,2]8.2.2}{[1][1][]1}}
\xrlabel{sec:8.2-absorption}{{8.2.3}{}{Absorption of a bosonic field}{Doc-Start}{}}
\xrlabel{sec:8.2-absorption@cref}{{[subsection][3][8,2]8.2.3}{[1][1][]1}}
\xrlabel{eq:8.2-abs}{{8.19}{}{Absorption of a bosonic field}{Doc-Start}{}}
\xrlabel{eq:8.2-abs@cref}{{[equation][19][8]8.19}{[1][1][]1}}
\xrlabel{sec:8.2-exothermic}{{8.2.4}{}{Negative energies and exothermic dark matter}{Doc-Start}{}}
\xrlabel{sec:8.2-exothermic@cref}{{[subsection][4][8,2]8.2.4}{[1][1][]1}}
\xrlabel{sec:8.2-cross}{{8.2.5}{}{Conditions for off-diagonal terms}{Doc-Start}{}}
\xrlabel{sec:8.2-cross@cref}{{[subsection][5][8,2]8.2.5}{[1][1][]1}}
\xrlabel{sec:8.3}{{8.3}{}{Exact relations}{Doc-Start}{}}
\xrlabel{sec:8.3@cref}{{[section][3][8]8.3}{[1][1][]1}}
\xrlabel{eq:8.3-bound}{{8.21}{}{Interference of two couplings}{Doc-Start}{}}
\xrlabel{eq:8.3-bound@cref}{{[equation][21][8]8.21}{[1][1][]1}}
\xrlabel{eq:8.3-moments}{{8.22}{}{Moment matrices}{Doc-Start}{}}
\xrlabel{eq:8.3-moments@cref}{{[equation][22][8]8.22}{[1][1][]1}}
\xrlabel{eq:8.3-M3}{{8.24}{}{The mixed third moment}{Doc-Start}{}}
\xrlabel{eq:8.3-M3@cref}{{[equation][24][8]8.24}{[1][1][]1}}
\xrlabel{eq:8.3-M3lim}{{8.25}{}{The mixed third moment}{Doc-Start}{}}
\xrlabel{eq:8.3-M3lim@cref}{{[equation][25][8]8.25}{[1][1][]1}}
\xrlabel{eq:8.3-M3sph}{{8.26}{}{The mixed third moment}{Doc-Start}{}}
\xrlabel{eq:8.3-M3sph@cref}{{[equation][26][8]8.26}{[1][1][]1}}
\xrlabel{we:8-third}{{8.4}{}{The mixed third moment}{Doc-Start}{}}
\xrlabel{we:8-third@cref}{{[workedexample][4][8]8.4}{[1][1][]1}}
\xrlabel{eq:8.3-mass}{{8.28}{}{Translation Ward identities}{Doc-Start}{}}
\xrlabel{eq:8.3-mass@cref}{{[equation][28][8]8.28}{[1][1][]1}}
\xrlabel{eq:8.3-soft}{{8.31}{}{Translation Ward identities}{Doc-Start}{}}
\xrlabel{eq:8.3-soft@cref}{{[equation][31][8]8.31}{[1][1][]1}}
\xrlabel{eq:8.3-boostward}{{8.32}{}{Translation Ward identities}{Doc-Start}{}}
\xrlabel{eq:8.3-boostward@cref}{{[equation][32][8]8.32}{[1][1][]1}}
\xrlabel{sec:8.4}{{8.4}{}{Screening for every coupling}{Doc-Start}{}}
\xrlabel{sec:8.4@cref}{{[section][4][8]8.4}{[1][1][]1}}
\xrlabel{sec:8.4-proper}{{8.4.1}{}{The proper response of a channel}{Doc-Start}{}}
\xrlabel{sec:8.4-proper@cref}{{[subsection][1][8,4]8.4.1}{[1][1][]1}}
\xrlabel{eq:8.4-screened}{{8.35}{}{The proper response of a channel}{Doc-Start}{}}
\xrlabel{eq:8.4-screened@cref}{{[equation][35][8]8.35}{[1][1][]1}}
\xrlabel{eq:8.4-rank1}{{8.36}{}{The proper response of a channel}{Doc-Start}{}}
\xrlabel{eq:8.4-rank1@cref}{{[equation][36][8]8.36}{[1][1][]1}}
\xrlabel{sec:8.4-couplings}{{8.4.2}{}{Four couplings}{Doc-Start}{}}
\xrlabel{sec:8.4-couplings@cref}{{[subsection][2][8,4]8.4.2}{[1][1][]1}}
\xrlabel{sec:8.4-exact}{{8.4.3}{}{Exactness of the loss function}{Doc-Start}{}}
\xrlabel{sec:8.4-exact@cref}{{[subsection][3][8,4]8.4.3}{[1][1][]1}}
\xrlabel{eq:8.4-exact}{{8.38}{}{Exactness of the loss function}{Doc-Start}{}}
\xrlabel{eq:8.4-exact@cref}{{[equation][38][8]8.38}{[1][1][]1}}
\xrlabel{we:8-nucleon}{{8.6}{}{Exactness of the loss function}{Doc-Start}{}}
\xrlabel{we:8-nucleon@cref}{{[workedexample][6][8]8.6}{[1][1][]1}}
\xrlabel{sec:8.4-lines}{{8.4.4}{}{Coulomb-line classes}{Doc-Start}{}}
\xrlabel{sec:8.4-lines@cref}{{[subsection][4][8,4]8.4.4}{[1][1][]1}}
\xrlabel{sec:8.4-mixing}{{8.4.5}{}{In-medium mixing}{Doc-Start}{}}
\xrlabel{sec:8.4-mixing@cref}{{[subsection][5][8,4]8.4.5}{[1][1][]1}}
\xrlabel{eq:8.4-zetaeff}{{8.39}{}{In-medium mixing}{Doc-Start}{}}
\xrlabel{eq:8.4-zetaeff@cref}{{[equation][39][8]8.39}{[1][1][]1}}
\xrlabel{sec:8.5}{{8.5}{}{The expansion hierarchy}{Doc-Start}{}}
\xrlabel{sec:8.5@cref}{{[section][5][8]8.5}{[1][1][]1}}
\xrlabel{eq:8.5-theta}{{8.40}{}{Small parameters}{Doc-Start}{}}
\xrlabel{eq:8.5-theta@cref}{{[equation][40][8]8.40}{[1][1][]1}}
\xrlabel{eq:8.5-xR}{{8.41}{}{Small parameters}{Doc-Start}{}}
\xrlabel{eq:8.5-xR@cref}{{[equation][41][8]8.41}{[1][1][]1}}
\xrlabel{eq:8.5-cumulants}{{8.42}{}{Lattice regimes}{Doc-Start}{}}
\xrlabel{eq:8.5-cumulants@cref}{{[equation][42][8]8.42}{[1][1][]1}}
\xrlabel{tab:8-params}{{8.4}{}{Small parameters of the expansion hierarchy, with their values for silicon in a detector at \SI {10}{mK}. Each controls one component of the truncation label $\mathfrak {T}$ (\cref {eq:8.5-label}). ``Introduced'' is the section where the physics first appears; ``Number'' is where it becomes a correction to a dark-matter rate. The values are estimates, with coefficients of order one. The inputs are $M = \SI {26.16}{GeV}$ (natural silicon) for $\kappa _{\mathrm {BO}}$, the optical hole vertex of \SI {0.25}{eV} and off-shell distances of 2 to \SI {2.5}{eV} at the absorption edge for $\theta _{\mathrm {ep}}$, the bond coordinate of \cref {sec:5.1} for $\eta _{\mathrm {anh}}$, the transverse acoustic phonon at X (\SI {18.57}{meV}) for $T/\omega _{\vect {Q}\nu }$, the density of states of \cref {sec:5.2} ($\langle \omega \rangle = \SI {40.9}{meV}$) and the natural isotope mixture for $x_R$, a kick of \SI {4e-21}{s} across the \SI {1.17}{eV} gap for $\Delta E\,\tau _{\mathrm {acc}}$, and an orbital of \SI {1}{\angstrom } for $\xi _D$}{Doc-Start}{}}
\xrlabel{tab:8-params@cref}{{[table][4][8]8.4}{[1][1][]1}}
\xrlabel{eq:8.5-width}{{8.43}{}{Lattice regimes}{Doc-Start}{}}
\xrlabel{eq:8.5-width@cref}{{[equation][43][8]8.43}{[1][1][]1}}
\xrlabel{eq:8.5-label}{{8.44}{}{The truncation label}{Doc-Start}{}}
\xrlabel{eq:8.5-label@cref}{{[equation][44][8]8.44}{[1][1][]1}}
\xrlabel{eq:8.5-base}{{8.45}{}{The truncation label}{Doc-Start}{}}
\xrlabel{eq:8.5-base@cref}{{[equation][45][8]8.45}{[1][1][]1}}
\xrlabel{we:8-label}{{8.7}{}{Many-body approximations as resummations}{Doc-Start}{}}
\xrlabel{we:8-label@cref}{{[workedexample][7][8]8.7}{[1][1][]1}}
\xrlabel{sec:8.6}{{8.6}{}{The decomposition theorem}{Doc-Start}{}}
\xrlabel{sec:8.6@cref}{{[section][6][8]8.6}{[1][1][]1}}
\xrlabel{eq:8.6-leak}{{8.47}{}{Three sectors of final states}{Doc-Start}{}}
\xrlabel{eq:8.6-leak@cref}{{[equation][47][8]8.47}{[1][1][]1}}
\xrlabel{eq:8.6-partition}{{8.48}{}{Three sectors of final states}{Doc-Start}{}}
\xrlabel{eq:8.6-partition@cref}{{[equation][48][8]8.48}{[1][1][]1}}
\xrlabel{eq:8.6-sector}{{8.49}{}{Three sectors of final states}{Doc-Start}{}}
\xrlabel{eq:8.6-sector@cref}{{[equation][49][8]8.49}{[1][1][]1}}
\xrlabel{eq:8.6-decomp}{{8.50}{}{Three sectors of final states}{Doc-Start}{}}
\xrlabel{eq:8.6-decomp@cref}{{[equation][50][8]8.50}{[1][1][]1}}
\xrlabel{eq:8.6-follow}{{8.51}{}{Densities in each sector}{Doc-Start}{}}
\xrlabel{eq:8.6-follow@cref}{{[equation][51][8]8.51}{[1][1][]1}}
\xrlabel{eq:8.6-electronic}{{8.52}{}{Densities in each sector}{Doc-Start}{}}
\xrlabel{eq:8.6-electronic@cref}{{[equation][52][8]8.52}{[1][1][]1}}
\xrlabel{eq:8.6-boost}{{8.53}{}{Densities in each sector}{Doc-Start}{}}
\xrlabel{eq:8.6-boost@cref}{{[equation][53][8]8.53}{[1][1][]1}}
\xrlabel{eq:8.6-fsum-e}{{8.54}{}{Partial sum rules}{Doc-Start}{}}
\xrlabel{eq:8.6-fsum-e@cref}{{[equation][54][8]8.54}{[1][1][]1}}
\xrlabel{eq:8.6-PMig}{{8.55}{}{Partial sum rules}{Doc-Start}{}}
\xrlabel{eq:8.6-PMig@cref}{{[equation][55][8]8.55}{[1][1][]1}}
\xrlabel{eq:8.6-fsum-k}{{8.56}{}{Partial sum rules}{Doc-Start}{}}
\xrlabel{eq:8.6-fsum-k@cref}{{[equation][56][8]8.56}{[1][1][]1}}
\xrlabel{eq:8.6-fsum-ek}{{8.57}{}{Partial sum rules}{Doc-Start}{}}
\xrlabel{eq:8.6-fsum-ek@cref}{{[equation][57][8]8.57}{[1][1][]1}}
\xrlabel{eq:8.6-amp}{{8.58}{}{Zero-phonon line and phonon sideband}{Doc-Start}{}}
\xrlabel{eq:8.6-amp@cref}{{[equation][58][8]8.58}{[1][1][]1}}
\xrlabel{eq:8.6-dyson}{{8.59}{}{Zero-phonon line and phonon sideband}{Doc-Start}{}}
\xrlabel{eq:8.6-dyson@cref}{{[equation][59][8]8.59}{[1][1][]1}}
\xrlabel{eq:8.6-ambiguity}{{8.61}{}{Zero-phonon line and phonon sideband}{Doc-Start}{}}
\xrlabel{eq:8.6-ambiguity@cref}{{[equation][61][8]8.61}{[1][1][]1}}
\xrlabel{eq:8.6-fano}{{8.62}{}{Zero-phonon line and phonon sideband}{Doc-Start}{}}
\xrlabel{eq:8.6-fano@cref}{{[equation][62][8]8.62}{[1][1][]1}}
\xrlabel{eq:8.6-bin}{{8.63}{}{Targets without a gap}{Doc-Start}{}}
\xrlabel{eq:8.6-bin@cref}{{[equation][63][8]8.63}{[1][1][]1}}
\xrlabel{we:8-sectors}{{8.8}{}{Targets without a gap}{Doc-Start}{}}
\xrlabel{we:8-sectors@cref}{{[workedexample][8][8]8.8}{[1][1][]1}}
\xrlabel{sec:8.7}{{8.7}{}{Final-state-resolved response}{Doc-Start}{}}
\xrlabel{sec:8.7@cref}{{[section][7][8]8.7}{[1][1][]1}}
\xrlabel{sec:8.7-classes}{{8.7.1}{}{Final-state classes}{Doc-Start}{}}
\xrlabel{sec:8.7-classes@cref}{{[subsection][1][8,7]8.7.1}{[1][1][]1}}
\xrlabel{eq:8.7-resolved}{{8.64}{}{Final-state classes}{Doc-Start}{}}
\xrlabel{eq:8.7-resolved@cref}{{[equation][64][8]8.64}{[1][1][]1}}
\xrlabel{eq:8.7-IP}{{8.66}{}{Pairs in an insulator}{Doc-Start}{}}
\xrlabel{eq:8.7-IP@cref}{{[equation][66][8]8.66}{[1][1][]1}}
\xrlabel{eq:8.7-factor}{{8.67}{}{Pairs in an insulator}{Doc-Start}{}}
\xrlabel{eq:8.7-factor@cref}{{[equation][67][8]8.67}{[1][1][]1}}
\xrlabel{eq:8.7-screened}{{8.68}{}{Pairs in an insulator}{Doc-Start}{}}
\xrlabel{eq:8.7-screened@cref}{{[equation][68][8]8.68}{[1][1][]1}}
\xrlabel{we:8-pairs}{{8.9}{}{Pairs in an insulator}{Doc-Start}{}}
\xrlabel{we:8-pairs@cref}{{[workedexample][9][8]8.9}{[1][1][]1}}
\xrlabel{sec:8.7-handoff}{{8.7.3}{}{Cascade probabilities of a final state}{Doc-Start}{}}
\xrlabel{sec:8.7-handoff@cref}{{[subsection][3][8,7]8.7.3}{[1][1][]1}}
\xrlabel{eq:8.7-handoff}{{8.69}{}{Cascade probabilities of a final state}{Doc-Start}{}}
\xrlabel{eq:8.7-handoff@cref}{{[equation][69][8]8.69}{[1][1][]1}}
\xrlabel{eq:8.7-cond}{{8.70}{}{Cascade probabilities of a final state}{Doc-Start}{}}
\xrlabel{eq:8.7-cond@cref}{{[equation][70][8]8.70}{[1][1][]1}}
\xrlabel{sec:8.8}{{8.8}{}{Uncertainty structure}{Doc-Start}{}}
\xrlabel{sec:8.8@cref}{{[section][8][8]8.8}{[1][1][]1}}
\xrlabel{eq:8.8-delta}{{8.71}{}{The uncertainty of a truncation}{Doc-Start}{}}
\xrlabel{eq:8.8-delta@cref}{{[equation][71][8]8.71}{[1][1][]1}}
\xrlabel{eq:8.8-M}{{8.72}{}{The uncertainty of a truncation}{Doc-Start}{}}
\xrlabel{eq:8.8-M@cref}{{[equation][72][8]8.72}{[1][1][]1}}
\xrlabel{eq:8.8-shift}{{8.74}{}{Scale, tilt and edge shifts}{Doc-Start}{}}
\xrlabel{eq:8.8-shift@cref}{{[equation][74][8]8.74}{[1][1][]1}}
\xrlabel{tab:8-nuisances}{{8.5}{}{Types of truncation uncertainty, with one truncation of each type in each sector of the decomposition \cref {eq:8.6-partition}: $\mathcal {P}_{\mathrm {ad}}$ (electrons following the nuclei; phonons and recoils), $\mathcal {P}_0$ (zero-phonon electronic excitations) and $\mathcal {P}_+$ (electronic excitations with phonons). The axis of the truncation label is given in parentheses. Sizes are for silicon, and the mark ``est.'' denotes an estimate from a small parameter of \cref {tab:8-params}. The residue entry in $\mathcal {P}_0$ and the finite-part entry in $\mathcal {P}_+$ carry one parameter, tied by \cref {eq:8.6-zsum}; the two $\mathcal {K}$ entries of $\mathcal {P}_{\mathrm {ad}}$, the two parts of one step of \cref {eq:5.6-multi}, also carry one parameter}{Doc-Start}{}}
\xrlabel{tab:8-nuisances@cref}{{[table][5][8]8.5}{[1][1][]1}}
\xrlabel{eq:8.8-fsum}{{8.76}{}{Exact relations as constraints}{Doc-Start}{}}
\xrlabel{eq:8.8-fsum@cref}{{[equation][76][8]8.76}{[1][1][]1}}
\xrlabel{eq:8.8-tilt}{{8.77}{}{Exact relations as constraints}{Doc-Start}{}}
\xrlabel{eq:8.8-tilt@cref}{{[equation][77][8]8.77}{[1][1][]1}}
\xrlabel{we:8-lf-nuisance}{{8.10}{}{Exact relations as constraints}{Doc-Start}{}}
\xrlabel{we:8-lf-nuisance@cref}{{[workedexample][10][8]8.10}{[1][1][]1}}
\xrlabel{app:C}{{C}{}{Mathematical tools}{Doc-Start}{}}
\xrlabel{app:C@cref}{{[appendix][3][2147483647]C}{[1][1][]1}}
\xrlabel{app:C.3}{{C.3}{}{Angular averages and spherical functions}{Doc-Start}{}}
\xrlabel{app:C.3@cref}{{[subappendix][3][2147483647,3]C.3}{[1][1][]1}}
\xrlabel{eq:C.3-addition}{{C.20}{}{Spherical harmonics}{Doc-Start}{}}
\xrlabel{eq:C.3-addition@cref}{{[equation][20][2147483647,3]C.20}{[1][1][]1}}
\xrlabel{eq:C.3-cubic}{{C.24}{}{Cubic harmonics}{Doc-Start}{}}
\xrlabel{eq:C.3-cubic@cref}{{[equation][24][2147483647,3]C.24}{[1][1][]1}}
\xrlabel{app:C.5}{{C.5}{}{Complex analysis for response functions}{Doc-Start}{}}
\xrlabel{app:C.5@cref}{{[subappendix][5][2147483647,3]C.5}{[1][1][]1}}
\xrlabel{eq:C.5-screening}{{C.38}{}{Direct and inverse dielectric functions}{Doc-Start}{}}
\xrlabel{eq:C.5-screening@cref}{{[equation][38][2147483647,3]C.38}{[1][1][]1}}
\xrlabel{eq:C.5-KKeps}{{C.39}{}{Direct and inverse dielectric functions}{Doc-Start}{}}
\xrlabel{eq:C.5-KKeps@cref}{{[equation][39][2147483647,3]C.39}{[1][1][]1}}
\xrlabel{app:C.7}{{C.7}{}{Matrix tools}{Doc-Start}{}}
\xrlabel{app:C.7@cref}{{[subappendix][7][2147483647,3]C.7}{[1][1][]1}}
\xrlabel{app:C.8}{{C.8}{}{Special functions}{Doc-Start}{}}
\xrlabel{app:C.8@cref}{{[subappendix][8][2147483647,3]C.8}{[1][1][]1}}
\xrlabel{eq:C.8-branch}{{C.73}{}{The Lindhard logarithms}{Doc-Start}{}}
\xrlabel{eq:C.8-branch@cref}{{[equation][73][2147483647,3]C.73}{[1][1][]1}}
\xrlabel{app:C.9}{{C.9}{}{Classical fields and static propagators}{Doc-Start}{}}
\xrlabel{app:C.9@cref}{{[subappendix][9][2147483647,3]C.9}{[1][1][]1}}
\xrlabel{eq:C.9-Gq}{{C.79}{}{The static Green's function}{Doc-Start}{}}
\xrlabel{eq:C.9-Gq@cref}{{[equation][79][2147483647,3]C.79}{[1][1][]1}}
\xrlabel{eq:C.9-mix}{{C.83}{}{Kinetic mixing}{Doc-Start}{}}
\xrlabel{eq:C.9-mix@cref}{{[equation][83][2147483647,3]C.83}{[1][1][]1}}
\xrlabel{app:C.10}{{C.10}{}{Moment problems}{Doc-Start}{}}
\xrlabel{app:C.10@cref}{{[subappendix][10][2147483647,3]C.10}{[1][1][]1}}
\xrlabel{eq:C.10-kernel}{{C.87}{}{Christoffel function and canonical representation}{Doc-Start}{}}
\xrlabel{eq:C.10-kernel@cref}{{[equation][87][2147483647,3]C.87}{[1][1][]1}}
\xrlabel{eq:C.10-rho}{{C.89}{}{Christoffel function and canonical representation}{Doc-Start}{}}
\xrlabel{eq:C.10-rho@cref}{{[equation][89][2147483647,3]C.89}{[1][1][]1}}
\xrlabel{eq:C.10-band}{{C.91}{}{The Markov--Krein band}{Doc-Start}{}}
\xrlabel{eq:C.10-band@cref}{{[equation][91][2147483647,3]C.91}{[1][1][]1}}
\xrlabel{eq:C.10-gauss}{{C.92}{}{The Markov--Krein band}{Doc-Start}{}}
\xrlabel{eq:C.10-gauss@cref}{{[equation][92][2147483647,3]C.92}{[1][1][]1}}
\xrlabel{eq:C.10-tail}{{C.93}{}{Tails of the band}{Doc-Start}{}}
\xrlabel{eq:C.10-tail@cref}{{[equation][93][2147483647,3]C.93}{[1][1][]1}}
\xrlabel{eq:C.10-stieltjes}{{C.94}{}{Distributions on the half line}{Doc-Start}{}}
\xrlabel{eq:C.10-stieltjes@cref}{{[equation][94][2147483647,3]C.94}{[1][1][]1}}
\xrlabel{app:C.11}{{C.11}{}{Generating functions, branching and renewal}{Doc-Start}{}}
\xrlabel{app:C.11@cref}{{[subappendix][11][2147483647,3]C.11}{[1][1][]1}}
\xrlabel{eq:C.11-thinfano}{{C.98}{}{Probability generating functions}{Doc-Start}{}}
\xrlabel{eq:C.11-thinfano@cref}{{[equation][98][2147483647,3]C.98}{[1][1][]1}}
\xrlabel{eq:C.11-functional}{{C.100}{}{Branching with types and deposits}{Doc-Start}{}}
\xrlabel{eq:C.11-functional@cref}{{[equation][100][2147483647,3]C.100}{[1][1][]1}}
\xrlabel{eq:C.11-mean}{{C.101}{}{Branching with types and deposits}{Doc-Start}{}}
\xrlabel{eq:C.11-mean@cref}{{[equation][101][2147483647,3]C.101}{[1][1][]1}}
\xrlabel{eq:C.11-second}{{C.102}{}{Branching with types and deposits}{Doc-Start}{}}
\xrlabel{eq:C.11-second@cref}{{[equation][102][2147483647,3]C.102}{[1][1][]1}}
\xrlabel{eq:C.11-key}{{C.104}{}{Renewal theory}{Doc-Start}{}}
\xrlabel{eq:C.11-key@cref}{{[equation][104][2147483647,3]C.104}{[1][1][]1}}
\xrlabel{eq:C.11-undershoot}{{C.105}{}{Renewal theory}{Doc-Start}{}}
\xrlabel{eq:C.11-undershoot@cref}{{[equation][105][2147483647,3]C.105}{[1][1][]1}}
\xrlabel{eq:C.11-linear}{{C.106}{}{Renewal theory}{Doc-Start}{}}
\xrlabel{eq:C.11-linear@cref}{{[equation][106][2147483647,3]C.106}{[1][1][]1}}
\xrlabel{app:C.12}{{C.12}{}{Linear programming and duality}{Doc-Start}{}}
\xrlabel{app:C.12@cref}{{[subappendix][12][2147483647,3]C.12}{[1][1][]1}}
\xrlabel{eq:C.12-sip}{{C.109}{}{Semi-infinite programs}{Doc-Start}{}}
\xrlabel{eq:C.12-sip@cref}{{[equation][109][2147483647,3]C.109}{[1][1][]1}}
\xrlabel{eq:C.12-certificate}{{C.110}{}{Semi-infinite programs}{Doc-Start}{}}
\xrlabel{eq:C.12-certificate@cref}{{[equation][110][2147483647,3]C.110}{[1][1][]1}}
\xrlabel{eq:C.12-mass}{{C.111}{}{Semi-infinite programs}{Doc-Start}{}}
\xrlabel{eq:C.12-mass@cref}{{[equation][111][2147483647,3]C.111}{[1][1][]1}}
\xrlabel{eq:C.12-onebin}{{C.112}{}{Extreme points and streams}{Doc-Start}{}}
\xrlabel{eq:C.12-onebin@cref}{{[equation][112][2147483647,3]C.112}{[1][1][]1}}
\xrlabel{eq:C.12-leak}{{C.113}{}{Discretization error}{Doc-Start}{}}
\xrlabel{eq:C.12-leak@cref}{{[equation][113][2147483647,3]C.113}{[1][1][]1}}
\xrlabel{app:C.13}{{C.13}{}{Interval orders}{Doc-Start}{}}
\xrlabel{app:C.13@cref}{{[subappendix][13][2147483647,3]C.13}{[1][1][]1}}
\xrlabel{eq:C.13-order}{{C.114}{}{Partial orders from intervals}{Doc-Start}{}}
\xrlabel{eq:C.13-order@cref}{{[equation][114][2147483647,3]C.114}{[1][1][]1}}
\xrlabel{eq:C.13-height}{{C.115}{}{Heights as tiers}{Doc-Start}{}}
\xrlabel{eq:C.13-height@cref}{{[equation][115][2147483647,3]C.115}{[1][1][]1}}
\xrlabel{app:C.14}{{C.14}{}{Likelihoods, intervals and profiling}{Doc-Start}{}}
\xrlabel{app:C.14@cref}{{[subappendix][14][2147483647,3]C.14}{[1][1][]1}}
\xrlabel{eq:C.14-poisson}{{C.116}{}{Poisson and Gaussian likelihoods}{Doc-Start}{}}
\xrlabel{eq:C.14-poisson@cref}{{[equation][116][2147483647,3]C.116}{[1][1][]1}}
\xrlabel{eq:C.14-ratio}{{C.118}{}{The likelihood ratio and intervals}{Doc-Start}{}}
\xrlabel{eq:C.14-ratio@cref}{{[equation][118][2147483647,3]C.118}{[1][1][]1}}
\xrlabel{eq:C.14-thresholds}{{C.119}{}{The likelihood ratio and intervals}{Doc-Start}{}}
\xrlabel{eq:C.14-thresholds@cref}{{[equation][119][2147483647,3]C.119}{[1][1][]1}}
\xrlabel{eq:C.14-profile}{{C.120}{}{Nuisance parameters and the profile likelihood}{Doc-Start}{}}
\xrlabel{eq:C.14-profile@cref}{{[equation][120][2147483647,3]C.120}{[1][1][]1}}
\xrlabel{eq:C.14-quadrature}{{C.121}{}{Nuisance parameters and the profile likelihood}{Doc-Start}{}}
\xrlabel{eq:C.14-quadrature@cref}{{[equation][121][2147483647,3]C.121}{[1][1][]1}}
\xrlabel{eq:C.14-pairwise}{{C.122}{}{Monte Carlo propagation}{Doc-Start}{}}
\xrlabel{eq:C.14-pairwise@cref}{{[equation][122][2147483647,3]C.122}{[1][1][]1}}
\xrlabel{app:C.15}{{C.15}{}{Dirac equation and Foldy--Wouthuysen reduction}{Doc-Start}{}}
\xrlabel{app:C.15@cref}{{[subappendix][15][2147483647,3]C.15}{[1][1][]1}}
\xrlabel{eq:C.15-HD}{{C.124}{}{The Dirac equation in a potential}{Doc-Start}{}}
\xrlabel{eq:C.15-HD@cref}{{[equation][124][2147483647,3]C.124}{[1][1][]1}}
\xrlabel{eq:C.15-spinor}{{C.125}{}{The Dirac equation in a potential}{Doc-Start}{}}
\xrlabel{eq:C.15-spinor@cref}{{[equation][125][2147483647,3]C.125}{[1][1][]1}}
\xrlabel{eq:C.15-bilinears}{{C.126}{}{Bilinears and their non-relativistic limits}{Doc-Start}{}}
\xrlabel{eq:C.15-bilinears@cref}{{[equation][126][2147483647,3]C.126}{[1][1][]1}}
\xrlabel{eq:C.15-darwin}{{C.131}{}{The contact density at a nucleus}{Doc-Start}{}}
\xrlabel{eq:C.15-darwin@cref}{{[equation][131][2147483647,3]C.131}{[1][1][]1}}
\xrlabel{eq:C.15-vertex}{{C.132}{}{The contact density at a nucleus}{Doc-Start}{}}
\xrlabel{eq:C.15-vertex@cref}{{[equation][132][2147483647,3]C.132}{[1][1][]1}}
\xrlabel{app:C.16}{{C.16}{}{The quantized radiation field}{Doc-Start}{}}
\xrlabel{app:C.16@cref}{{[subappendix][16][2147483647,3]C.16}{[1][1][]1}}
\xrlabel{eq:C.16-soft}{{C.138}{}{The one-photon emission rate}{Doc-Start}{}}
\xrlabel{eq:C.16-soft@cref}{{[equation][138][2147483647,3]C.138}{[1][1][]1}}
\xrlabel{eq:C.16-medium}{{C.139}{}{A photon mode in a medium}{Doc-Start}{}}
\xrlabel{eq:C.16-medium@cref}{{[equation][139][2147483647,3]C.139}{[1][1][]1}}
\xrlabel{app:C.17}{{C.17}{}{Point and space groups}{Doc-Start}{}}
\xrlabel{app:C.17@cref}{{[subappendix][17][2147483647,3]C.17}{[1][1][]1}}
\xrlabel{eq:C.17-orth}{{C.141}{}{Representations and characters}{Doc-Start}{}}
\xrlabel{eq:C.17-orth@cref}{{[equation][141][2147483647,3]C.141}{[1][1][]1}}
\xrlabel{eq:C.17-ninv}{{C.142}{}{Representations and characters}{Doc-Start}{}}
\xrlabel{eq:C.17-ninv@cref}{{[equation][142][2147483647,3]C.142}{[1][1][]1}}
\xrlabel{eq:C.17-polar}{{C.143}{}{Polar and axial tensors}{Doc-Start}{}}
\xrlabel{eq:C.17-polar@cref}{{[equation][143][2147483647,3]C.143}{[1][1][]1}}
\xrlabel{eq:C.17-gyration}{{C.144}{}{Responses odd in q}{Doc-Start}{}}
\xrlabel{eq:C.17-gyration@cref}{{[equation][144][2147483647,3]C.144}{[1][1][]1}}
\xrlabel{we:C-gyration}{{C.1}{}{Responses odd in q}{Doc-Start}{}}
\xrlabel{we:C-gyration@cref}{{[workedexample][1][2147483647,3]C.1}{[1][1][]1}}
\xrlabel{eq:C.17-glide}{{C.145}{}{The diamond space group}{Doc-Start}{}}
\xrlabel{eq:C.17-glide@cref}{{[equation][145][2147483647,3]C.145}{[1][1][]1}}
\xrlabel{eq:C.17-extinct}{{C.146}{}{The diamond space group}{Doc-Start}{}}
\xrlabel{eq:C.17-extinct@cref}{{[equation][146][2147483647,3]C.146}{[1][1][]1}}
\xrlabel{eq:C.17-parity}{{C.147}{}{The diamond space group}{Doc-Start}{}}
\xrlabel{eq:C.17-parity@cref}{{[equation][147][2147483647,3]C.147}{[1][1][]1}}
\xrlabel{eq:C.17-selection}{{C.148}{}{The diamond space group}{Doc-Start}{}}
\xrlabel{eq:C.17-selection@cref}{{[equation][148][2147483647,3]C.148}{[1][1][]1}}
\xrlabel{app:C.18}{{C.18}{}{Continuum elasticity}{Doc-Start}{}}
\xrlabel{app:C.18@cref}{{[subappendix][18][2147483647,3]C.18}{[1][1][]1}}
\xrlabel{eq:C.18-bulk}{{C.153}{}{Strain, stress and the elastic tensor}{Doc-Start}{}}
\xrlabel{eq:C.18-bulk@cref}{{[equation][153][2147483647,3]C.153}{[1][1][]1}}
\xrlabel{eq:C.18-christoffel}{{C.154}{}{The Christoffel problem}{Doc-Start}{}}
\xrlabel{eq:C.18-christoffel@cref}{{[equation][154][2147483647,3]C.154}{[1][1][]1}}
\xrlabel{eq:C.18-cubic}{{C.155}{}{The Christoffel problem}{Doc-Start}{}}
\xrlabel{eq:C.18-cubic@cref}{{[equation][155][2147483647,3]C.155}{[1][1][]1}}
\xrlabel{eq:C.18-rayleigh}{{C.156}{}{The Christoffel problem}{Doc-Start}{}}
\xrlabel{eq:C.18-rayleigh@cref}{{[equation][156][2147483647,3]C.156}{[1][1][]1}}
\xrlabel{eq:C.18-transmission}{{C.159}{}{Acoustic impedance and reflection}{Doc-Start}{}}
\xrlabel{eq:C.18-transmission@cref}{{[equation][159][2147483647,3]C.159}{[1][1][]1}}
\xrlabel{eq:C.18-snell}{{C.160}{}{Acoustic impedance and reflection}{Doc-Start}{}}
\xrlabel{eq:C.18-snell@cref}{{[equation][160][2147483647,3]C.160}{[1][1][]1}}
\xrlabel{app:C.19}{{C.19}{}{Magnetic neutron scattering}{Doc-Start}{}}
\xrlabel{app:C.19@cref}{{[subappendix][19][2147483647,3]C.19}{[1][1][]1}}
\xrlabel{eq:C.19-field}{{C.161}{}{The magnetic cross-section}{Doc-Start}{}}
\xrlabel{eq:C.19-field@cref}{{[equation][161][2147483647,3]C.161}{[1][1][]1}}
\xrlabel{eq:C.19-amplitude}{{C.162}{}{The magnetic cross-section}{Doc-Start}{}}
\xrlabel{eq:C.19-amplitude@cref}{{[equation][162][2147483647,3]C.162}{[1][1][]1}}
\xrlabel{eq:C.19-cross}{{C.163}{}{The magnetic cross-section}{Doc-Start}{}}
\xrlabel{eq:C.19-cross@cref}{{[equation][163][2147483647,3]C.163}{[1][1][]1}}
\xrlabel{eq:C.19-spinonly}{{C.164}{}{Ions and the dipole form factor}{Doc-Start}{}}
\xrlabel{eq:C.19-spinonly@cref}{{[equation][164][2147483647,3]C.164}{[1][1][]1}}
\xrlabel{eq:C.19-dipole}{{C.165}{}{Ions and the dipole form factor}{Doc-Start}{}}
\xrlabel{eq:C.19-dipole@cref}{{[equation][165][2147483647,3]C.165}{[1][1][]1}}
\xrlabel{eq:C.19-polar}{{C.166}{}{Polarization analysis}{Doc-Start}{}}
\xrlabel{eq:C.19-polar@cref}{{[equation][166][2147483647,3]C.166}{[1][1][]1}}
\xrlabel{app:C.20}{{C.20}{}{Dissipative two-level systems and kinks}{Doc-Start}{}}
\xrlabel{app:C.20@cref}{{[subappendix][20][2147483647,3]C.20}{[1][1][]1}}
\xrlabel{eq:C.20-H}{{C.167}{}{The spin--boson model}{Doc-Start}{}}
\xrlabel{eq:C.20-H@cref}{{[equation][167][2147483647,3]C.167}{[1][1][]1}}
\xrlabel{eq:C.20-J}{{C.168}{}{The spin--boson model}{Doc-Start}{}}
\xrlabel{eq:C.20-J@cref}{{[equation][168][2147483647,3]C.168}{[1][1][]1}}
\xrlabel{eq:C.20-ohmic}{{C.169}{}{The spin--boson model}{Doc-Start}{}}
\xrlabel{eq:C.20-ohmic@cref}{{[equation][169][2147483647,3]C.169}{[1][1][]1}}
\xrlabel{eq:C.20-rate}{{C.172}{}{Incoherent tunneling rate}{Doc-Start}{}}
\xrlabel{eq:C.20-rate@cref}{{[equation][172][2147483647,3]C.172}{[1][1][]1}}
\xrlabel{eq:C.20-ohmrate}{{C.173}{}{Incoherent tunneling rate}{Doc-Start}{}}
\xrlabel{eq:C.20-ohmrate@cref}{{[equation][173][2147483647,3]C.173}{[1][1][]1}}
\xrlabel{eq:C.20-Dr}{{C.174}{}{Incoherent tunneling rate}{Doc-Start}{}}
\xrlabel{eq:C.20-Dr@cref}{{[equation][174][2147483647,3]C.174}{[1][1][]1}}
\xrlabel{eq:C.20-FK}{{C.175}{}{The Frenkel--Kontorova chain}{Doc-Start}{}}
\xrlabel{eq:C.20-FK@cref}{{[equation][175][2147483647,3]C.175}{[1][1][]1}}
\xrlabel{eq:C.20-kink}{{C.178}{}{The Frenkel--Kontorova chain}{Doc-Start}{}}
\xrlabel{eq:C.20-kink@cref}{{[equation][178][2147483647,3]C.178}{[1][1][]1}}
\xrlabel{eq:C.20-Ekink}{{C.179}{}{The Frenkel--Kontorova chain}{Doc-Start}{}}
\xrlabel{eq:C.20-Ekink@cref}{{[equation][179][2147483647,3]C.179}{[1][1][]1}}
\xrlabel{eq:C.20-mass}{{C.180}{}{The Frenkel--Kontorova chain}{Doc-Start}{}}
\xrlabel{eq:C.20-mass@cref}{{[equation][180][2147483647,3]C.180}{[1][1][]1}}
\xrlabel{eq:C.20-EPN}{{C.182}{}{The Frenkel--Kontorova chain}{Doc-Start}{}}
\xrlabel{eq:C.20-EPN@cref}{{[equation][182][2147483647,3]C.182}{[1][1][]1}}
\xrlabel{app:D}{{D}{}{Long derivations}{Doc-Start}{}}
\xrlabel{app:D@cref}{{[appendix][4][2147483647]D}{[1][1][]1}}
\xrlabel{app:D.1}{{D.1}{}{Standard-halo integrals}{Doc-Start}{}}
\xrlabel{app:D.1@cref}{{[subappendix][1][2147483647,4]D.1}{[1][1][]1}}
\xrlabel{eq:D.1-eta1}{{D.2}{}{The function \texorpdfstring {$\eta (\vmin )$}{eta(vmin)}}{Doc-Start}{}}
\xrlabel{eq:D.1-eta1@cref}{{[equation][2][2147483647,4]D.2}{[1][1][]1}}
\xrlabel{eq:D.1-eta2}{{D.3}{}{The function \texorpdfstring {$\eta (\vmin )$}{eta(vmin)}}{Doc-Start}{}}
\xrlabel{eq:D.1-eta2@cref}{{[equation][3][2147483647,4]D.3}{[1][1][]1}}
\xrlabel{app:D.6}{{D.4}{}{The f-sum rule}{Doc-Start}{}}
\xrlabel{app:D.6@cref}{{[subappendix][4][2147483647,4]D.4}{[1][1][]1}}
\xrlabel{app:D.9}{{D.7}{}{The Mermin relaxation-time dielectric function}{Doc-Start}{}}
\xrlabel{app:D.9@cref}{{[subappendix][7][2147483647,4]D.7}{[1][1][]1}}
\xrlabel{app:D.12}{{D.10}{}{The Bethe--Salpeter equation from the ladder}{Doc-Start}{}}
\xrlabel{app:D.12@cref}{{[subappendix][10][2147483647,4]D.10}{[1][1][]1}}
\xrlabel{eq:D.12-blocks}{{D.45}{}{The transition-basis Hamiltonian}{Doc-Start}{}}
\xrlabel{eq:D.12-blocks@cref}{{[equation][45][2147483647,4]D.45}{[1][1][]1}}
\xrlabel{app:D.14}{{D.12}{}{The Fr\"ohlich vertex}{Doc-Start}{}}
\xrlabel{app:D.14@cref}{{[subappendix][12][2147483647,4]D.12}{[1][1][]1}}
\xrlabel{eq:D.14-gL}{{D.51}{}{The Fr\"ohlich vertex}{Doc-Start}{}}
\xrlabel{eq:D.14-gL@cref}{{[equation][51][2147483647,4]D.51}{[1][1][]1}}
\xrlabel{app:D.20}{{D.14}{}{BCS gap and coherence factors}{Doc-Start}{}}
\xrlabel{app:D.20@cref}{{[subappendix][14][2147483647,4]D.14}{[1][1][]1}}
\xrlabel{app:D.21}{{D.15}{}{Conserving BCS density response}{Doc-Start}{}}
\xrlabel{app:D.21@cref}{{[subappendix][15][2147483647,4]D.15}{[1][1][]1}}
\xrlabel{app:D.22}{{D.16}{}{Polarization of Dirac electrons}{Doc-Start}{}}
\xrlabel{app:D.22@cref}{{[subappendix][16][2147483647,4]D.16}{[1][1][]1}}
\xrlabel{app:D.23}{{D.17}{}{Effective-mass theory of shallow donors}{Doc-Start}{}}
\xrlabel{app:D.23@cref}{{[subappendix][17][2147483647,4]D.17}{[1][1][]1}}
\xrlabel{eq:D.23-vo}{{D.63}{}{Valleys and the central cell}{Doc-Start}{}}
\xrlabel{eq:D.23-vo@cref}{{[equation][63][2147483647,4]D.63}{[1][1][]1}}
\xrlabel{app:D.24}{{D.18}{}{The operator basis}{Doc-Start}{}}
\xrlabel{app:D.24@cref}{{[subappendix][18][2147483647,4]D.18}{[1][1][]1}}
\xrlabel{app:D.25}{{D.19}{}{Master rate formulas at finite temperature}{Doc-Start}{}}
\xrlabel{app:D.25@cref}{{[subappendix][19][2147483647,4]D.19}{[1][1][]1}}
\xrlabel{app:D.27}{{D.21}{}{Screening for every external coupling}{Doc-Start}{}}
\xrlabel{app:D.27@cref}{{[subappendix][21][2147483647,4]D.21}{[1][1][]1}}
\xrlabel{app:D.28}{{D.22}{}{The decomposition theorem}{Doc-Start}{}}
\xrlabel{app:D.28@cref}{{[subappendix][22][2147483647,4]D.22}{[1][1][]1}}
\xrlabel{eq:D.28-dM}{{D.86}{}{The electron-following projector}{Doc-Start}{}}
\xrlabel{eq:D.28-dM@cref}{{[equation][86][2147483647,4]D.86}{[1][1][]1}}
\xrlabel{eq:D.28-boost}{{D.87}{}{The kick and the electronic sectors}{Doc-Start}{}}
\xrlabel{eq:D.28-boost@cref}{{[equation][87][2147483647,4]D.87}{[1][1][]1}}
\xrlabel{chap:23}{{23}{}{Density-functional theory}{Doc-Start}{}}
\xrlabel{chap:23@cref}{{[chapter][23][]23}{[1][1][]1}}
\xrlabel{app:D.5}{{23.2}{}{Hohenberg--Kohn theorems and the Kohn--Sham construction}{Doc-Start}{}}
\xrlabel{app:D.5@cref}{{[section][2][23]23.2}{[1][1][]1}}
\xrlabel{app:E.1}{{23.3}{}{Kohn--Sham calculations in practice}{Doc-Start}{}}
\xrlabel{app:E.1@cref}{{[section][3][23]23.3}{[1][1][]1}}
\xrlabel{tab:E-functionals}{{23.1}{}{Lattice constants and band gaps of silicon, germanium and gallium arsenide for three functionals, without spin--orbit coupling. Lattice constants from LDA and PBE are all-electron values from Ref.~\cite [Table~II]{csonka2009assessing}, those from HSE from Ref.~\cite [Table~I]{lucero2012improved}. Gaps are evaluated at each functional's own relaxed lattice constant, from Ref.~\cite [Table~I]{lucero2012improved}, whose LDA lattice constant of silicon is \SI {5.410}{\angstrom }; a zero gap means that the functional makes the crystal a metal. The static-lattice constants are the measured ones extrapolated to zero temperature with the zero-point anharmonic expansion removed~\cite [Table~II]{csonka2009assessing}; they are not measured at any temperature, and they lie below the lattice constants measured at low temperature. They are used only inside calculations, as the cell of a static lattice. The room-temperature lattice constants are measured: silicon~\cite [lattice parameter of silicon]{mohr2025codata}, germanium~\cite [Sect.~II\,A]{rideau2006strained} and GaAs~\cite {vurgaftman2001band}; they are the values for observables and the cell of the shared lattice setting of \cref {app:E.17}. The measured gaps are low-temperature values~\cite [Table~I]{lucero2012improved}}{Doc-Start}{}}
\xrlabel{tab:E-functionals@cref}{{[table][1][23]23.1}{[1][1][]1}}
\xrlabel{app:E.3}{{23.4}{}{Quasiparticle corrections}{Doc-Start}{}}
\xrlabel{app:E.3@cref}{{[section][4][23]23.4}{[1][1][]1}}
\xrlabel{sec:3.5}{{24.1}{}{Plane waves and the core region}{Doc-Start}{}}
\xrlabel{sec:3.5@cref}{{[section][1][24]24.1}{[1][1][]1}}
\xrlabel{eq:3.5-rhoAE}{{24.6}{}{Projector augmented waves}{Doc-Start}{}}
\xrlabel{eq:3.5-rhoAE@cref}{{[equation][6][24]24.6}{[1][1][]1}}
\xrlabel{fig:3-core}{{24.2}{}{How the electrons of each atomic shell are distributed in momentum, for (a)~silicon and (b)~germanium: electrons per keV of momentum magnitude, $N_{nl}\,p^2|\phi _{nl}(p)|^2$, with $N_{nl}$ the occupation of the shell and $\phi _{nl}$ the radial momentum-space orbital normalized to $\int \dd p\,p^2|\phi _{nl}|^2 = 1$. The orbitals are those of the free atoms, from an all-electron, scalar-relativistic calculation with the PBE functional (configurations $[\mathrm {Ne}]\,3s^23p^2$ and $[\mathrm {Ar}]\,3d^{10}4s^24p^2$; germanium $1s$, which opens above $q = \SI {4}{MeV}$, is not shown). Black: the valence shell. Dotted vertical lines: the smallest momentum transfer at which dark matter of any mass can lift an electron of the shell of the same color into the conduction band, $(\Eg + E_{\mathrm {core}})/v_{\mathrm {max}}$ (twice that at the threshold mass), with $v_{\mathrm {max}} = 2.650\times 10^{-3}$, gaps of 1.17 and \SI {0.744}{eV} and measured binding energies relative to the valence-band maximum~\cite [Table~1-1]{williams2009electron}. Above about \SI {5}{keV} the core electrons outnumber the valence electrons in both atoms. The silicon $2p$ shell opens at \SI {38}{keV}, in the tail of its distribution; the germanium $3d$ shell opens at \SI {11}{keV}, near its peak}{Doc-Start}{}}
\xrlabel{fig:3-core@cref}{{[figure][2][24]24.2}{[1][1][]1}}
\xrlabel{eq:3.5-HBP}{{24.8}{}{Relativistic corrections}{Doc-Start}{}}
\xrlabel{eq:3.5-HBP@cref}{{[equation][8][24]24.8}{[1][1][]1}}
\xrlabel{we:3-Zalpha}{{24.2}{}{Relativistic corrections}{Doc-Start}{}}
\xrlabel{we:3-Zalpha@cref}{{[workedexample][2][24]24.2}{[1][1][]1}}
\xrlabel{app:E.2}{{24.2}{}{Construction and testing of pseudopotentials}{Doc-Start}{}}
\xrlabel{app:E.2@cref}{{[section][2][24]24.2}{[1][1][]1}}
\xrlabel{chap:25}{{25}{}{Phonons from first principles}{Doc-Start}{}}
\xrlabel{chap:25@cref}{{[chapter][25][]25}{[1][1][]1}}
\xrlabel{sec:5.4}{{25.1}{}{Frozen phonons and linear response}{Doc-Start}{}}
\xrlabel{sec:5.4@cref}{{[section][1][25]25.1}{[1][1][]1}}
\xrlabel{app:E.4}{{25.2}{}{Density-functional perturbation theory in practice}{Doc-Start}{}}
\xrlabel{app:E.4@cref}{{[section][2][25]25.2}{[1][1][]1}}
\xrlabel{eq:E.4-drho}{{25.19}{}{Wavefunction changes beyond the density}{Doc-Start}{}}
\xrlabel{eq:E.4-drho@cref}{{[equation][19][25]25.19}{[1][1][]1}}
\xrlabel{app:E.8}{{25.3}{}{Anharmonic and finite-temperature methods}{Doc-Start}{}}
\xrlabel{app:E.8@cref}{{[section][3][25]25.3}{[1][1][]1}}
\xrlabel{eq:E.8-wl}{{25.29}{}{Special displacements}{Doc-Start}{}}
\xrlabel{eq:E.8-wl@cref}{{[equation][29][25]25.29}{[1][1][]1}}
\xrlabel{we:E-zpr}{{25.3}{}{Zero-point renormalization from special displacements}{Doc-Start}{}}
\xrlabel{we:E-zpr@cref}{{[workedexample][3][25]25.3}{[1][1][]1}}
\xrlabel{app:E.6}{{26.1}{}{Dielectric response from first principles}{Doc-Start}{}}
\xrlabel{app:E.6@cref}{{[section][1][26]26.1}{[1][1][]1}}
\xrlabel{eq:E.6-L}{{26.3}{}{Computed quantities}{Doc-Start}{}}
\xrlabel{eq:E.6-L@cref}{{[equation][3][26]26.3}{[1][1][]1}}
\xrlabel{app:E.7}{{26.2}{}{Crystal form factors}{Doc-Start}{}}
\xrlabel{app:E.7@cref}{{[section][2][26]26.2}{[1][1][]1}}
\xrlabel{eq:E.7-vNL}{{26.9}{}{Velocity matrix elements and nonlocal potentials}{Doc-Start}{}}
\xrlabel{eq:E.7-vNL@cref}{{[equation][9][26]26.9}{[1][1][]1}}
\xrlabel{eq:E.7-core}{{26.11}{}{Core and semicore states}{Doc-Start}{}}
\xrlabel{eq:E.7-core@cref}{{[equation][11][26]26.11}{[1][1][]1}}
\xrlabel{sec:3.6}{{27.1}{}{Wannier functions}{Doc-Start}{}}
\xrlabel{sec:3.6@cref}{{[section][1][27]27.1}{[1][1][]1}}
\xrlabel{app:E.5}{{27.2}{}{Wannier functions in practice}{Doc-Start}{}}
\xrlabel{app:E.5@cref}{{[section][2][27]27.2}{[1][1][]1}}
\xrlabel{app:E.9}{{28.1}{}{Magnetic and superconducting parameters}{Doc-Start}{}}
\xrlabel{app:E.9@cref}{{[section][1][28]28.1}{[1][1][]1}}
\xrlabel{app:E.10}{{28.2}{}{Defects}{Doc-Start}{}}
\xrlabel{app:E.10@cref}{{[section][2][28]28.2}{[1][1][]1}}
\xrlabel{eq:E.10-ohm}{{28.15}{}{Tunneling defects in metals}{Doc-Start}{}}
\xrlabel{eq:E.10-ohm@cref}{{[equation][15][28]28.15}{[1][1][]1}}
\xrlabel{app:E.11}{{29.1}{}{Recoil molecular dynamics}{Doc-Start}{}}
\xrlabel{app:E.11@cref}{{[section][1][29]29.1}{[1][1][]1}}
\xrlabel{eq:E.11-Est}{{29.6}{}{Stored energy and the quench protocol}{Doc-Start}{}}
\xrlabel{eq:E.11-Est@cref}{{[equation][6][29]29.6}{[1][1][]1}}
\xrlabel{tab:E-md-si}{{29.1}{}{Recoil molecular dynamics in silicon: the inputs and the quantities a production study converges}{Doc-Start}{}}
\xrlabel{tab:E-md-si@cref}{{[table][1][29]29.1}{[1][1][]1}}
\xrlabel{app:E.12}{{29.2}{}{Phonon and charge transport Monte Carlo}{Doc-Start}{}}
\xrlabel{app:E.12@cref}{{[section][2][29]29.2}{[1][1][]1}}
\xrlabel{eq:E.12-iso}{{29.8}{}{Phonon rates}{Doc-Start}{}}
\xrlabel{eq:E.12-iso@cref}{{[equation][8][29]29.8}{[1][1][]1}}
\xrlabel{eq:E.12-anh}{{29.9}{}{Phonon rates}{Doc-Start}{}}
\xrlabel{eq:E.12-anh@cref}{{[equation][9][29]29.9}{[1][1][]1}}
\xrlabel{app:E.13}{{29.3}{}{Hot-carrier cascades from first principles}{Doc-Start}{}}
\xrlabel{app:E.13@cref}{{[section][3][29]29.3}{[1][1][]1}}
\xrlabel{eq:E.13-gph}{{29.12}{}{The electron--phonon kernel}{Doc-Start}{}}
\xrlabel{eq:E.13-gph@cref}{{[equation][12][29]29.12}{[1][1][]1}}
\xrlabel{eq:E.13-ab}{{29.13}{}{The single-parameter model}{Doc-Start}{}}
\xrlabel{eq:E.13-ab@cref}{{[equation][13][29]29.13}{[1][1][]1}}
\xrlabel{app:E.17}{{30.1}{}{Lattice libraries}{Doc-Start}{}}
\xrlabel{app:E.17@cref}{{[section][1][30]30.1}{[1][1][]1}}
\xrlabel{app:E.19}{{30.2}{}{Anharmonic, disordered and quantum solids}{Doc-Start}{}}
\xrlabel{app:E.19@cref}{{[section][2][30]30.2}{[1][1][]1}}
\xrlabel{app:E.20}{{30.3}{}{Magnetic, superconducting and Dirac targets}{Doc-Start}{}}
\xrlabel{app:E.20@cref}{{[section][3][30]30.3}{[1][1][]1}}
\xrlabel{app:E.20-magnet}{{30.3.4}{}{Magnet inputs from DFT~\texorpdfstring {$+\,U$}{+U}}{Doc-Start}{}}
\xrlabel{app:E.20-magnet@cref}{{[subsection][4][30,3]30.3.4}{[1][1][]1}}
\xrlabel{app:E.21}{{30.4}{}{Dopants, color centers and small targets}{Doc-Start}{}}
\xrlabel{app:E.21@cref}{{[section][4][30]30.4}{[1][1][]1}}
\xrlabel{app:E.22}{{30.5}{}{Defects and two-level systems}{Doc-Start}{}}
\xrlabel{app:E.22@cref}{{[section][5][30]30.5}{[1][1][]1}}
\xrlabel{eq:E.22-SHR}{{30.5}{}{Release paths of eV-scale defects}{Doc-Start}{}}
\xrlabel{eq:E.22-SHR@cref}{{[equation][5][30]30.5}{[1][1][]1}}
\xrlabel{eq:E.22-Spol}{{30.6}{}{Crowdions and dumbbells}{Doc-Start}{}}
\xrlabel{eq:E.22-Spol@cref}{{[equation][6][30]30.6}{[1][1][]1}}
\xrlabel{app:E.14}{{31.1}{}{Phonon-assisted electron recoil from first principles}{Doc-Start}{}}
\xrlabel{app:E.14@cref}{{[section][1][31]31.1}{[1][1][]1}}
\xrlabel{app:E.15}{{31.2}{}{Excitons and core states}{Doc-Start}{}}
\xrlabel{app:E.15@cref}{{[section][2][31]31.2}{[1][1][]1}}
\xrlabel{eq:E.15-disp}{{31.2}{}{Exciton dispersion and the hydrogenic series}{Doc-Start}{}}
\xrlabel{eq:E.15-disp@cref}{{[equation][2][31]31.2}{[1][1][]1}}
\xrlabel{eq:E.15-ste}{{31.3}{}{Self-trapped excitons}{Doc-Start}{}}
\xrlabel{eq:E.15-ste@cref}{{[equation][3][31]31.3}{[1][1][]1}}
\xrlabel{app:E.16}{{31.3}{}{Electronic response libraries}{Doc-Start}{}}
\xrlabel{app:E.16@cref}{{[section][3][31]31.3}{[1][1][]1}}
\xrlabel{eq:E.16-dLF}{{31.4}{}{Local fields in the phonon-assisted window}{Doc-Start}{}}
\xrlabel{eq:E.16-dLF@cref}{{[equation][4][31]31.4}{[1][1][]1}}
\xrlabel{eq:E.16-step}{{31.5}{}{Static screening with lattice relaxation}{Doc-Start}{}}
\xrlabel{eq:E.16-step@cref}{{[equation][5][31]31.5}{[1][1][]1}}
\xrlabel{app:E.18}{{31.4}{}{Migdal and nonadiabatic quantities}{Doc-Start}{}}
\xrlabel{app:E.18@cref}{{[section][4][31]31.4}{[1][1][]1}}
\xrlabel{eq:E.18-X}{{31.9}{}{Couplings not proportional to mass}{Doc-Start}{}}
\xrlabel{eq:E.18-X@cref}{{[equation][9][31]31.9}{[1][1][]1}}
\xrlabel{app:E.26}{{31.5}{}{Coherent and second-order absorption}{Doc-Start}{}}
\xrlabel{app:E.26@cref}{{[section][5][31]31.5}{[1][1][]1}}
\xrlabel{app:E.23}{{32.1}{}{Nuclear-recoil partition and charge yield}{Doc-Start}{}}
\xrlabel{app:E.23@cref}{{[section][1][32]32.1}{[1][1][]1}}
\xrlabel{eq:E.24-dft}{{32.3}{}{A library of pair-number distributions}{Doc-Start}{}}
\xrlabel{eq:E.24-dft@cref}{{[equation][3][32]32.3}{[1][1][]1}}
\xrlabel{app:E.25}{{32.3}{}{Event-level transport}{Doc-Start}{}}
\xrlabel{app:E.25@cref}{{[section][3][32]32.3}{[1][1][]1}}
\xrlabel{app:E.27}{{32.4}{}{Halo, background and nuclear-spin inputs}{Doc-Start}{}}
\xrlabel{app:E.27@cref}{{[section][4][32]32.4}{[1][1][]1}}
\xrlabel{chap:33}{{33}{}{Libraries and open calculations}{Doc-Start}{}}
\xrlabel{chap:33@cref}{{[chapter][33][]33}{[1][1][]1}}
\xrlabel{sec:21.6}{{33.1}{}{Shared response libraries}{Doc-Start}{}}
\xrlabel{sec:21.6@cref}{{[section][1][33]33.1}{[1][1][]1}}
\xrlabel{tab:22-calc}{{33.2}{}{Open calculations that would change a printed result, limit or ranking of the book, ordered by the result each would change; the remaining calculation steps are stated in the sections indexed by \cref {tab:22-gaps}. Calculation: what is to be computed, with the section that specifies it. Result: the number or statement of the book that the calculation replaces, with its present size; masses are dark-matter masses. Check: the exact relation (sum rule, positivity, symmetry zero, conservation law) and the control against data that the result must pass. Section: where the item is stated. The guide is the ranking of crystals by figure of merit in \cref {sec:21.7}, its survivors are the eleven crystals with a documented detector record, and $P(\mathrm {first})$ is the probability that a crystal ranks first in a row of the guide}{Doc-Start}{}}
\xrlabel{tab:22-calc@cref}{{[table][2][33]33.2}{[1][1][]1}}
